\documentclass[11pt]{article}
\usepackage[margin=1in]{geometry}
\usepackage[T1]{fontenc}\usepackage[utf8]{inputenc}\usepackage{lmodern}
\usepackage{amsmath,amssymb,amsthm,mathtools,booktabs,microtype}
\usepackage[colorlinks=true,linkcolor=blue,urlcolor=blue,citecolor=blue]{hyperref}
\hypersetup{pdftitle={Shallow comparison oracles for deep Bayesian networks},pdfauthor={ChatGPT-6 Astra; Claude Opus 5.5; Kaarel Hänni; Dmitry Vaintrob}}
\newtheorem{theorem}{Theorem}[section]\newtheorem{lemma}[theorem]{Lemma}\newtheorem{proposition}[theorem]{Proposition}\newtheorem{corollary}[theorem]{Corollary}
\newcommand{\E}{\mathbb E}\newcommand{\R}{\mathbb R}\newcommand{\KL}{\operatorname{KL}}\newcommand{\op}{\mathrm{op}}
\title{Shallow comparison oracles for\\deep Bayesian networks}
\author{ChatGPT-6 Astra \and Claude Opus 5.5 \and Kaarel H\"anni \and Dmitry Vaintrob\thanks{Authors are listed in alphabetical order.}}
\date{17 November 2026}
\begin{document}\maketitle
\begin{abstract}
We study one finite dense tanh network with independent weight coordinates, a calibrated Gaussian likelihood, and raw-output squared loss. The prior may be Gaussian, a fixed smooth strongly log-concave coordinate law with bounded curvature, or a fixed finite Gaussian-mixture coordinate law with common positive component variance, under the same variance scalings. At every learner depth, its sample complexity for competing with a shallow comparison network is polynomial in comparison width. At canonical mean field, the depth factor is $\exp(O(\sqrt H))$; at variance prefactor $\nu\le1$ it is $\exp(O(\sqrt H+(1-\nu)H))$. The polynomial degree and constants fix input dimension, comparison coefficient cap, accuracy, and confidence. The guarantee holds at every learner width at least as large as the comparison width, for arbitrary bounded IID labels on the unit sphere. A corollary permits a fixed sub-Gaussian label envelope, including Gaussian noise, without changing the observations or likelihood. The bound has leading coefficient one on the best comparison risk and includes the ordinary raw posterior mean. No output-envelope assumption is imposed on the comparison functions, and the prior has no robustness restriction or mixture over architectures. The proof combines Gaussian Jacobian moments, a covariance construction that transports a weak feature through many layers, and finite-parameter covers at smaller widths. Same-sign attenuation of large comparison outputs is enough for squared-loss comparison. A deterministic choice of the comparison-width cap from sample count gives a single calibrated learner competing with width $\operatorname{poly}(n)$ shallow networks at fixed dimension and depth. For non-Gaussian coordinates, an exact-baseline matrix moment estimate and a quantitative conditional central limit argument preserve the same depth and width bounds. The Gaussian-mixture extension includes genuinely non-log-concave priors. These are statistical guarantees for exact posteriors, not efficient posterior-sampling algorithms.
\end{abstract}

\section{Model and main result}
Increasing learner depth need not preserve a Bayesian approximation guarantee. Even a shallow target can become difficult when the prior's native signals contract through many layers. We give a quantitative upper bound for this situation. The teacher has one hidden layer, while the learner may have arbitrarily many; these depths are not identified. We first prove the Gaussian case, then extend the same guarantee to smooth non-Gaussian independent coordinate priors in Theorem~\ref{thm:nongaussian-shallow-oracle}.

Let inputs lie on the unit sphere $\mathbb S^{d-1}\subset[-1,1]^d$, with $\bar x=(1,x)$, and let labels lie in $[-1,1]$. Their joint distribution $D$ is arbitrary and the training sample is IID. Write
\[
 R(f)=\E_D(f(X)-Y)^2,\qquad
 \widehat R_S(f)=n^{-1}\sum_{i=1}^n(f(X_i)-Y_i)^2.
\]
The width-$N$, hidden-depth-$H$ learner is
\[
 h^1=\tanh(W^1\bar x),\quad
 h^l=\tanh(W^lh^{l-1})\ (2\le l\le H),\quad
 z_w=v^\top h^H.
\]
There are no hidden bias parameters beyond the first-layer constant coordinate. For $0<\nu\le1$, the independent centered Gaussian prior variances are
\[
 \nu/d\text{ in }W^1,\qquad \nu/N\text{ internally},\qquad
 \nu/N^2\text{ in }v.
\]
Put $\omega=N/\nu$, $s=1/\omega=\nu/N$, and $\tau=\sqrt\nu$. The posterior is
\begin{equation}
 d\Pi_{H,N,\nu,S}(w)=Z^{-1}
       e^{-\omega\widehat R_S(z_w)/(2\lambda)}d\mu_{H,N,\nu}(w).
 \label{eq:posterior}
\end{equation}
This is a Gaussian likelihood with variance $n\lambda/\omega$. Both likelihood and prediction use the raw real-valued output. The ordinary posterior mean is $\bar z(x)=\int z_w(x)d\Pi(w)$.

Let $\mathcal G_{C,B}$ be the shallow class
\[
 g(x)=\sum_{j=1}^C c_j\tanh(a_j^\top\bar x),
 \qquad |c_j|,|a_{jk}|\le B.
\]
There is no additional output-envelope restriction: the automatic bound is $|g|\le CB$. Zero padding makes these classes nested in $C$.

\begin{theorem}[An all-depth shallow comparison oracle]
\label{thm:main-shallow}
Fix $d\ge1$, $B\ge1$, and $0<\epsilon,\delta<1$. There are constants $a,c>0$ and $0<\ell_-,\ell_0,\ell_+<1$, depending only on these fixed resources, with the following property. For every comparison width $C\ge1$ and every hidden depth $H\ge1$, put
\[
 D_C=1+\log(2+C),\qquad
 N_0=\left\lceil(2+C)^a(H+1)^a\right\rceil,
 \qquad
 T_{H,C,\nu}=1+(1-\nu)^2H+D_C^2/H,
\]
and $j_{H,C,\nu}=\lceil\log_2 T_{H,C,\nu}\rceil$. Use the likelihood calibration
\[
 \lambda_{H,N,\nu}=
 \begin{cases}
 \ell_- (2+C)^{-a}(H+1)^{-a},&N<N_0,\\
 \ell_0[H\exp(aD_C/H)]^{-1},&N\ge N_0,\ H<aD_C,\\
 \ell_+2^{-j_{H,C,\nu}},&N\ge N_0,\ H\ge aD_C.
 \end{cases}
\]
For every sample count $n$, one event of probability at least $1-\delta$ gives
\begin{equation}
 \int R(z_w)d\Pi_{H,N,\nu,S}
 \le\inf_{g\in\mathcal G_{C,B}}R(g)+\epsilon
 \label{eq:main-shallow}
\end{equation}
simultaneously for every $N\ge C$ and every $0<\nu\le1$ satisfying
\begin{equation}
 n\ge(2+C)^c\exp\{c[\sqrt H+(1-\nu)H]\}.
 \label{eq:main-shallow-budget}
\end{equation}
The same inequality holds for $R(\bar z)$.
\end{theorem}

At $\nu=1$, the sample bound is polynomial in comparison width times $\exp(O(\sqrt H))$. The depth factor stays of this order when $1-\nu=O(H^{-1/2})$. For any fixed smaller prefactor it is exponential in $H$. A power-law scaling $\nu=N^{1-\gamma}$ with $\gamma\ge1$ is a special case, although the theorem permits arbitrary prefactors. The polynomial degrees and constants may be large and depend on dimension; we do not claim joint polynomial dependence on all teacher resources.

\begin{corollary}[A sample-count calibration without knowing teacher width]
\label{cor:sample-cap}
Fix a learner architecture and prefactor, and let $c$ be as in Theorem~\ref{thm:main-shallow}. If the following set is nonempty, define
\[
 C_n=\max\{1\le C\le N:
          (2+C)^c e^{c[\sqrt H+(1-\nu)H]}\le n\}.
\]
Use the theorem's single likelihood calibration with $C=C_n$. Its risk is at most
$\inf_{g\in\mathcal G_{C_n,B}}R(g)+\epsilon$ with probability at least $1-\delta$. At fixed $d,H,B,\epsilon,\delta,\nu$, the eligible comparison width grows as a positive power of sample count until it reaches $N$.
\end{corollary}

The cap is a deterministic function of sample count and known model quantities, chosen before observing the data. This corollary introduces no prior over comparison widths: there is one fixed Gaussian prior on the learner's weights and one calibrated likelihood. The comparison networks have one hidden layer; the statement does not assert the same width exponent for arbitrary-depth comparison MLPs.

\section{Basic variational and Gaussian estimates}
The finite posterior obeys the exact variational identity
\begin{equation}
 F_S(\lambda):=-\frac{2\lambda}{\omega}\log Z
 =\inf_\rho\left\{\int\widehat R_S(z)d\rho
                     +\frac{2\lambda}{\omega}\KL(\rho\Vert\mu)\right\}.
 \label{eq:variational}
\end{equation}
In particular posterior empirical risk is at most $F_S(\lambda)$. Bounded tanh features and the independent Gaussian readout give
\[
 \E_\mu R(z)\le2,\qquad \E_\mu R(z)^2\le24,
 \qquad Z\ge e^{-\omega/\lambda}.
\]
The first two follow by conditioning on hidden features, where the output is centered Gaussian with variance at most $\nu/N\le1$; Jensen gives the evidence bound.

Define total operator energy
\[
 Q_{\rm op}=\|W^1\|_F^2/N+
             \sum_{l=2}^H\|W^l\|_\op^2+N\|v\|^2.
\]
For absolute $c_0,C_0>0$,
\begin{equation}
 \E_\mu e^{c_0\omega Q_{\rm op}}\le e^{C_0N(H+1)}.
 \label{eq:relu-total-op-mgf}
\end{equation}
The first and last factors are chi-square moment generating functions. For each internal standardized Gaussian matrix $G$, the Gaussian tail
$\Pr(\|G\|_\op>2\sqrt N+t)\le e^{-t^2/2}$ gives
$\E e^{c_0\|G\|_\op^2}\le e^{C_0N}$ for small $c_0$. Its scaling cancels against $\omega=N/\nu$. Independence multiplies the factors. These estimates hold at every finite width.

\section{A Gaussian gradient bound sensitive to comparison cost}
\label{sec:sobolev-depth}
A small relative-entropy comparison cost is useful only with a statistical bound that retains its depth dependence. We obtain such a bound from the input Jacobian. It uses moments of the actual input Jacobian, followed by a Sobolev embedding, rather than the product of full matrix operator norms.

\begin{lemma}[Gaussian input-gradient moments]
\label{lem:gaussian-gradient}
For every globally one-Lipschitz hidden activation, every even integer $p\ge2$, and Lebesgue-almost every input $x$,
\begin{equation}
 \E_\mu\|\nabla z_w(x)\|_2^p
 \le\left(\frac p\omega\right)^{p/2}
                     \exp\left(\frac{Hp^2}{2\omega}\right).
 \label{eq:gradient-moment}
\end{equation}
No condition on the activation's value at zero is needed.
\end{lemma}
\begin{proof}
First take a continuously differentiable activation and write $J_l$ for the $N$-by-$d$ hidden Jacobian. Its diagonal activation derivative has operator norm at most one, so
$\|J_1\|_F\le\|W^1_{\mathrm{input}}\|_F$ and
$\|J_l\|_F\le\|W^lJ_{l-1}\|_F$.
There is no independence assumption on the activation derivative.

For deterministic $J$ and a standard $N$-by-$N$ Gaussian matrix $G$, convexity over the normalized eigenvalues of $J^\top J$ gives
\[
 \E\|GJ\|_F^p\le\E\chi_N^p\,\|J\|_F^p.
\]
Indeed the squared norm is a weighted sum of independent chi-square variables with $N$ degrees of freedom; Jensen bounds its $p/2$ moment by the rank-one case. For even $p$,
$\E\chi_N^p=\prod_{j=0}^{p/2-1}(N+2j)\le(N+p)^{p/2}$.
For $0<t\le1$ and $r\ge0$, elementary calculus gives
$t(1+r)\le e^{tr}$.
It follows that
\[
 \tau^p\E(\chi_N/\sqrt N)^p
 \le[\tau^2(1+p/N)]^{p/2}
 \le e^{p^2/(2\omega)}.
\]
Conditioning successively on preceding layers, and using the analogous $\chi_{Nd}$ first-layer estimate, yields
$\E\|J_H\|_F^p\le N^{p/2}e^{Hp^2/(2\omega)}$.
The independent readout satisfies
\[
 \E_v\|v^\top J_H\|_2^p
 \le(\tau/N)^p\E|G_1|^p\,\|J_H\|_F^p
 \le(\tau\sqrt p/N)^p\|J_H\|_F^p.
\]
Combining proves \eqref{eq:gradient-moment}. For a general Lipschitz activation, apply the argument to smooth Lipschitz mollifications. They converge uniformly; the corresponding networks converge uniformly on each compact input domain for each fixed parameter vector. Weak lower semicontinuity of the Sobolev seminorm and Fatou's lemma give the same integrated moment estimate, which is the form used below. Equivalently the weak chain rule gives the pointwise almost-everywhere statement.
\end{proof}

\begin{lemma}[A common H\"older class outside exponentially small prior mass]
\label{lem:holder-prior}
Assume the hidden activation has range in $[-1,1]$. Fix $K\ge1$. If
$\omega\sqrt{K/H}\ge\max\{d,1\}$, then outside a prior set of mass at most $2e^{-\omega K}$,
\[
 \|z\|_\infty\le\sqrt{4K+2\log2},\qquad
 |z(x)-z(y)|\le C_d e^{C\sqrt{HK}}\|x-y\|_2^{1/2}.
\]
The constants and function class are independent of $N$ and $\tau$.
\end{lemma}
\begin{proof}
Use normalized Lebesgue measure on the cube and integrate \eqref{eq:gradient-moment}. Put $a=\omega\sqrt{K/H}$ and $p=2\lceil a\rceil$, so $p\ge2d$ and $p\le4a$. Markov's inequality gives
\[
 \Pr_\mu\left\{\|\nabla z\|_{L^p}>
 \sqrt{p/\omega}\exp\left[\frac{Hp}{2\omega}
                              +\frac{\omega K}{p}\right]\right\}
 \le e^{-\omega K}.
\]
The displayed cutoff is at most $C e^{C\sqrt{HK}}$. Morrey's inequality on a cube has a dimension-only constant uniformly for $p\ge2d$. It gives H\"older exponent $1-d/p\ge1/2$; weakening to exponent $1/2$ changes only a dimension constant. The continuous Sobolev representative is the actual network function.

For the amplitude, bounded hidden features give $\|z\|_\infty^2\le N\|v\|^2$. The exact Gaussian identity
$\E e^{\omega N\|v\|^2/4}=2^{N/2}$ and $N\le\omega$ imply
\[
 \Pr_\mu\{N\|v\|^2>4K+2\log2\}\le e^{-\omega K}.
\]
Intersect the events.
\end{proof}

\begin{proposition}[Generalization from a small normalized comparison cost]
\label{prop:sobolev-oracle}
Fix $d$, $0<\epsilon,\delta<1$, a bounded one-Lipschitz hidden activation, and a likelihood scale $0<\lambda\le1$. Set
$K=C[1+\lambda^{-1}+\log(1/\epsilon)]$ with a sufficiently large absolute $C$.
There is a common sample event on which every finite $N$ and every $0<\tau\le1$ satisfy
\begin{equation}
 \int R(z_w)d\Pi\le
 -\frac{2\lambda}{\omega}\log Z+\epsilon,
 \label{eq:sobolev-certificate}
\end{equation}
provided
\begin{equation}
 n\ge \operatorname{poly}_d(H,K,\epsilon^{-1},\log(1/\delta))
       +e^{C_d\sqrt{HK}}
          \operatorname{poly}_d(K,\epsilon^{-1},\log(1/\delta)).
 \label{eq:sobolev-budget}
\end{equation}
The polynomials have fixed degrees, independent of depth. The event has probability at least $1-\delta$.

Consequently, for a comparison function with a fixed output envelope, a uniformly accurate comparison law with
$2\KL(\nu\Vert\mu)/\omega\le A$ yields a leading-one oracle inequality at fixed accuracy with sample bound
$\operatorname{poly}_d(H,A)+\exp(C_d\sqrt{H(A+1)})$, after choosing $\lambda$ of order $1/(A+1)$. Its constants may also depend on that fixed output envelope. This consequence applies only at widths where the stated comparison exists. Any separate width threshold and its complementary branch must also be included.
\end{proposition}
\begin{proof}
For the large-effective-width case use Lemma~\ref{lem:holder-prior}. Its amplitude is $M=O(\sqrt{K+1})$. A spatial grid gives a uniform output cover at accuracy $b$ with
\[
 \log|\mathcal D|\le
 \left(\frac{C_d e^{C\sqrt{HK}}}{b}\right)^{2d}
                         \log(CM/b).
\]
Take $b=c\epsilon/(M+1)$. Finite-class Hoeffding for squared loss then gives the desired good-set loss deviation at the exponential term in \eqref{eq:sobolev-budget}.

If $\omega\sqrt{K/H}<\max\{d,1\}$, then
$N\le\omega<\max\{d,1\}\sqrt{H/K}$.
There are only polynomially many possible widths. The Gaussian moment estimate \eqref{eq:relu-total-op-mgf} gives a total-operator-energy cutoff $Q_{\mathrm{op}}\le C(H+K)$ with prior failure at most $e^{-\omega K}$; also impose the same readout amplitude cutoff. The coordinate magnitudes are at most $\sqrt{CN(H+K)}$, the parameter count is $O(dN+HN^2)$, and the logarithm of a parameter-to-output Lipschitz constant is $O(H\log(N+H+K+2))$. Finite parameter grids for these widths therefore have logarithmic size polynomial in $H,d,K,\log(1/b)$. Clipping the dictionary outputs to $[-M,M]$ preserves approximation accuracy. This gives the polynomial term in \eqref{eq:sobolev-budget}.

Both dictionaries are fixed independently of the sample, width, and scale. Their concentration event is common to every learner. At any fixed learner, the auxiliary prior-good set has complement at most $C e^{-\omega K}$. Since bounded hidden features imply $\E_\mu R(z)^2\le24$, and Jensen gives $Z\ge e^{-\omega/\lambda}$, Cauchy--Schwarz bounds the posterior raw-risk contribution of that complement by
\[
 \sqrt{24C}\exp[\omega/\lambda-\omega K/2]\le\epsilon/4
\]
with the stated choice of $K$ (enlarging the fixed constant if needed).
Integrate the good-set deviation and add this tail. Finally
$\int\widehat R\,d\Pi\le-2\lambda\log Z/\omega$ by the posterior variational identity. This proves \eqref{eq:sobolev-certificate}.

For the oracle consequence, use a comparison law whose approximation makes its empirical loss at most that of the comparison function plus a prescribed fraction of the requested accuracy. The variational identity gives the certificate at most that empirical loss plus $\lambda A$. Choose $\lambda=c_\epsilon/(A+1)$, and apply scalar concentration to the fixed bounded comparison function. Then $K=O_\epsilon(A+1)$, giving the stated budget. The ordinary raw posterior mean obeys the same bound by Jensen.
\end{proof}

\begin{corollary}[Bounded observation of general Lipschitz hidden layers]
\label{cor:sobolev-bounded-observation}
Proposition~\ref{prop:sobolev-oracle} also holds for any fixed globally one-Lipschitz hidden activation, whether bounded or centered, if a one-Lipschitz map into $[-1,1]$ is applied to the final output in both the likelihood and prediction. The certificate and risks then refer to those bounded outputs. Constants may depend on the hidden activation's value at zero.
\end{corollary}
\begin{proof}
The derivative moment in Lemma~\ref{lem:gaussian-gradient} does not require bounded hidden activations, and the final link cannot increase it. The output envelope is now one. For the small-effective-width branch, operator-energy coordinate caps still give parameter-to-output Lipschitz constants with logarithms polynomial in depth; the fixed activation offset adds only the corresponding constant to the layerwise growth estimate. The final link bounds every loss by four, so $Z\ge e^{-2\omega/\lambda}$ and a prior bad set of mass $Ce^{-\omega K}$ contributes at most $4C e^{2\omega/\lambda-\omega K}$ to posterior risk. Enlarge the fixed coefficient in $K$. All covering and variational arguments are unchanged.
\end{proof}

\section{Scalar facts for the critical tanh variance chain}
\label{sec:echo-scalar-preliminaries}
Write
\[
 F(q)=\E\tanh^2(\sqrt q Z),\qquad Z\sim N(0,1),\qquad
 \chi(q)=\E\tanh'(\sqrt q Z)=1-F(q).
\]
The following estimates also make the constants in the covariance schedules explicit in their dependence on a fixed input dimension.

\begin{lemma}[Variance drift and derivative products]
\label{lem:echo-scalar-facts}
The function $F$ is continuous and strictly increasing on $(0,\infty)$, with $0\le F'(q)\le1$. For absolute $c,c'>0$ and $0<q\le1$,
\[
 c'q\le F(q),\qquad q-2q^2\le F(q)\le q-cq^2.
\]
Fix $q_0\in(0,1)$ and put $q_{j+1}=F(q_j)$. There are positive constants depending only on $q_0$ such that
\[
 \frac c{j+1}\le q_j\le\frac C{j+1},\qquad
 c\sqrt{q_m/q_0}\le\prod_{j=0}^{m-1}(1-q_{j+1})
                \le C\sqrt{q_m/q_0}.
\]
In particular the product is comparable to $(m+1)^{-1/2}$.
\end{lemma}
\begin{proof}
Strict monotonicity and continuity follow pointwise from $\tanh^2(\sqrt q z)$ and bounded convergence. The heat-semigroup differentiation identity, justified by bounded derivatives, gives
\[
 F'(q)=\tfrac12\E(\tanh^2)''(\sqrt q Z).
\]
With $u=\tanh^2 t$, the half second derivative is $1-4u+3u^2\le1$ for $0\le u\le1$. Monotonicity supplies the lower derivative bound.

For $0\le t\le1$, concavity gives $\tanh t\ge t/2$, and hence
\[
 t-\tanh t=\int_0^t\tanh^2 s\,ds\ge t^3/12.
\]
Thus $t^2-\tanh^2t\ge t^4/12$. Restricting the Gaussian integral to $|Z|\le1$ proves the upper drift bound and the lower linear bound. For the other drift inequality use
$\tanh^2t\ge t^2-(2/3)t^4$ for all real $t$. If its right side is positive, $|t|\le\sqrt{3/2}$ and the inequality follows by squaring $\tanh t\ge t-t^3/3$ for $t\ge0$; this last bound follows by integrating $\tanh'(t)=1-\tanh^2t\ge1-t^2$. Where the right side is nonpositive the claim is immediate. Gaussian averaging gives $F(q)\ge q-2q^2$.

The reciprocal increment is bounded above and below by positive constants:
\[
 c\le\frac1{q_{j+1}}-\frac1{q_j}
 =\frac{q_j-F(q_j)}{q_jF(q_j)}\le2/c'.
\]
Summation proves the bounds on $q_j$.

Taylor expansion with a bounded remainder gives
\[
 F(q)=q-2q^2+\tfrac{17}{3}q^3+O(q^4).
\]
For example expand $\tanh^2t=t^2-(2/3)t^4+(17/45)t^6+O(t^8)$ with a globally valid bounded-remainder constant and average the Gaussian moments. Consequently
\[
 J(q):=\frac{q[1-F(q)]^2}{F(q)}
       =1-\tfrac23q^2+O(q^3).
\]
Every $J(q_j)$ is positive. Since $\sum_jq_j^2<\infty$, the product $\prod_{j<m}J(q_j)$ is bounded above and below by positive constants independent of $m$: its tail logarithm is absolutely summable and its finite initial segment is positive. The exact telescoping identity
\[
 \prod_{j=0}^{m-1}(1-q_{j+1})^2
   =\frac{q_m}{q_0}\prod_{j=0}^{m-1}J(q_j)
\]
proves the derivative-product estimate.
\end{proof}

\section{A depth-independent comparison for a fixed shallow teacher}
\label{sec:continuous-echo}
This section uses the unchanged canonical Gaussian prior: first-layer variance $1/d$, internal variance $1/N$, and readout variance $1/N^2$. Inputs lie on the unit sphere in $\mathbb R^d$, and the first layer includes its usual constant coordinate. The conclusion is a pointwise mean-square comparison; it does not require a simultaneous approximation event on the whole sphere.

\begin{theorem}[A fixed teacher at arbitrarily large depth]
\label{thm:continuous-echo}
Fix $d$, a finite shallow tanh function
\[
 g(x)=\sum_{j=1}^C c_j\tanh(a_j^{\mathsf T}(1,x)),
\]
and an accuracy $\varepsilon>0$. There exist finite constants $A,H_0,K,p$, depending on these fixed objects but not on learner depth or width, such that for every $H\ge H_0$ and every $N\ge K H^p$ there is a comparison law $\rho\ll\mu_{H,N}$ satisfying
\[
 \frac2N\KL(\rho\|\mu_{H,N})\le A,
 \qquad
 \sup_{x\in\mathbb S^{d-1}}\E_\rho|z_w(x)-g(x)|^2\le\varepsilon^2.
\]
The constants can be chosen uniformly over all teachers of a fixed width and a fixed coefficient bound. The same law therefore approximates $g$ in expected squared loss under every input distribution on the sphere and on every deterministic training sample. No restriction is imposed on the learner's prior support. The first-layer marginal of the comparison remains exactly the native Gaussian law.
\end{theorem}

\subsection{A reverse schedule with an almost linear response}
Let $F(q)=\E\tanh^2(\sqrt q Z)$, put $q_{-1}=2/d$ and $q_0=F(q_{-1})$, and recursively set $q_j=F(q_{j-1})$. Constants in this section may depend on the fixed dimension. The critical drift gives
\[
 \frac{c_d}{j+1}\le q_j\le\frac{C_d}{j+1},\qquad
 q_j-q_{j+1}\le C_dq_j^2.
\]
For $H=2M+1$, use $M$ native transitions followed by $M$ covariance spikes. The reverse steps, indexed by their remaining count $t=M,M-1,\ldots,1$, have
\[
 R_t^2=\frac{q_{t-2}}{q_t},\qquad \alpha_t=R_t^2-1\ge0.
\]
Thus their reference variances run from $q_M$ back to $q_0$. If $H$ is even, append one native transition to the odd-depth schedule with $H-1$ hidden layers. The fixed final native transition changes none of the estimates below. We have
\[
 \alpha_t\le\frac{C_d}{t+1},\qquad
 \sum_t\{R_t^2-1-\log R_t^2\}\le C_d,
 \qquad
 \prod_{t=1}^M R_t
 =\left(\frac{q_{-1}q_0}{q_{M-1}q_M}\right)^{1/2}
 \asymp_d H.
\]
The first bound and the convergent square sum imply the entropy estimate; the finitely many final factors are bounded constants.

For zero initial correlation the hidden variances continue to follow the native sequence, even during the spike schedule. Define $b_\ell=q_{\ell-1}$ at hidden layer $\ell$ and let
\[
 L_1=1,\qquad L_{\ell+1}=R_\ell(1-b_{\ell+1})L_\ell,
\]
where $R_\ell=1$ on native transitions and is the displayed reverse factor on a spike transition. The small-variance identity
\[
 \frac{q(1-F(q))^2}{F(q)}=1+O(q^2)
\]
and summability of $q_j^2$ show that
\[
 \prod_{\ell=1}^{H-1}(1-b_{\ell+1})\asymp_d H^{-1/2},
 \qquad L_H\asymp_d\sqrt H.
\]
The constants are strictly positive because every factor is positive and only finitely many initial variances lie outside a fixed small interval. For a state immediately before a reverse step with $t$ steps remaining, telescoping the remaining $R$ factors and bounding the remaining native derivative product above and below gives
\begin{equation}
 \frac{L_\ell}{L_H}\le\frac{C_d}{t+1}.
 \label{eq:continuous-response-ratio}
\end{equation}
Indeed, all native variances in the reverse half are $O(1/H)$, so their derivative product over at most $H$ transitions is bounded below by a positive dimension-dependent constant. The remaining spike product is comparable to $t+1$.

Consider the exact scalar recursion for a prescribed initial correlation $r_1$:
\[
 q_1=q_0,\quad
 v_\ell=q_\ell+(R_\ell^2-1)r_\ell^2,\quad
 q_{\ell+1}=F(v_\ell),\quad
 r_{\ell+1}=R_\ell r_\ell(1-q_{\ell+1}).
\]
Choose $r_1=s/L_H$ for a signed intended final signal $s$. During the native burn $q_\ell=b_\ell$. Throughout the schedule, $q_\ell\ge b_\ell$ and $|r_\ell|\le |s|L_\ell/L_H$. Since $0\le F'\le1$, the excess variance $\Delta_\ell=q_\ell-b_\ell$ satisfies
\[
 0\le\Delta_{\ell+1}\le\Delta_\ell+
             (R_\ell^2-1)s^2(L_\ell/L_H)^2.
\]
Each perturbation contributes to at most $t+1$ subsequent excess variances. Equation~\eqref{eq:continuous-response-ratio} therefore gives
\[
 \sum_{\ell=1}^H\Delta_\ell
 \le C_ds^2\sum_{t\ge1}\frac{t+1}{(t+1)^3}
 \le C_ds^2.
\]
Here $F'\le1$ follows by differentiating the Gaussian expectation: $F'(q)=\frac12\E(\tanh^2)''(\sqrt qZ)\le1$, while monotonicity gives the lower bound. Dividing the exact correlation product by its linearization yields
\[
 \frac{r_H}{s}=\prod_{\ell=2}^H
       \left(1-\frac{\Delta_\ell}{1-b_\ell}\right).
\]
All factors lie in $[0,1]$. Since $1-b_\ell\ge1-q_0>0$, the elementary product inequality gives the uniform response estimate
\begin{equation}
 |r_H-s|\le C_d|s|^3.
 \label{eq:continuous-cubic-response}
\end{equation}
For $s=0$ the assertion follows directly. The scalar recursion is admissible for every sufficiently large $H$ whenever $s$ lies in a fixed bounded interval, because $r_1^2\le q_0$.

\subsection{Native Gaussian features encode the small initial signal}
Let $W\sim N(0,I_{d+1}/d)$. There is a bounded measurable function $\psi(W)$ such that
\[
 g_0(x)=\E[\psi(W)\tanh(W^{\mathsf T}(1,x))]
 \quad\hbox{satisfies}\quad
 \sup_{x\in\mathbb S^{d-1}}|g_0(x)-g(x)|\le\varepsilon/4.
\]
To construct it explicitly, choose small balls $B_j$ around the finitely many $a_j$, put $p_j=\Pr(W\in B_j)>0$, and set $\psi=\sum_j c_j\mathbf1_{B_j}/p_j$. The $1$-Lipschitz property of tanh and $\|(1,x)\|=\sqrt2$ make the approximation uniform as the ball radii decrease. In particular $M_\psi=\|\psi\|_\infty<\infty$ is fixed independently of $H$. For a fixed teacher width and coefficient bound, use a common ball radius proportional to $\varepsilon/(1+\sum_j|c_j|)$. The Gaussian masses of these balls have a uniform positive lower bound because their centers lie in a fixed compact set. Thus $M_\psi$ and every later comparison constant are uniform over that teacher class.

Choose a fixed readout factor $\mathcal B\ge1+\|g\|_\infty+\|g_0\|_\infty$ so large that
\[
 C_d\|g_0\|_\infty^3/\mathcal B^2\le\varepsilon/4.
\]
The intended scalar signal at input $x$ is $s(x)=g_0(x)/\mathcal B$. Draw the first-layer rows $W_i$ from their unchanged Gaussian law and put
\[
 a_i=\frac{\psi(W_i)}{\mathcal B L_H\sqrt N}.
\]
For all sufficiently large $H$, $\|a\|\le M_\psi/(\mathcal B L_H)\le1/2$ deterministically. Conditional on the first-layer weights, take an auxiliary random unit vector $v$ uniform on the sphere perpendicular to $a$ and set
\[
 e_1=a+\sqrt{1-\|a\|^2}\,v.
\]
If $a=0$, use a uniform unit vector in the whole space. This gives an exact unit direction. At any fixed input, Hoeffding concentration for the bounded feature average and the elementary spherical-coordinate tail give
\[
 \left|\frac{h^1(x)^{\mathsf T}e_1}{\sqrt N}
           -\frac{g_0(x)}{\mathcal B L_H}\right|\le\delta,
 \qquad
 \left|\frac{\|h^1(x)\|^2}{N}-q_0\right|\le\delta
\]
except on probability at most $C\exp(-cN\delta^2)$, with constants depending only on the fixed teacher, accuracy, and dimension. These bounds are pointwise in $x$ but their constants are uniform over the sphere. They do not require a simultaneous event for every input.

\subsection{Finite-width propagation and entropy}
At an internal transition, conditionally on the past and the current unit direction $e$, change each row covariance from $I/N$ to
\[
 N^{-1}\{I+(R_\ell^2-1)ee^{\mathsf T}\}.
\]
Let $\xi_i$ be the standardized Gaussian component of row $i$ in direction $e$, and take the next direction to be $\xi/\|\xi\|$. Gaussian integration by parts gives exactly the scalar recursion above for the conditional expectations before the empirical normalization of $\xi$. Conditional concentration of the bounded squares, the sub-Gaussian products $\xi_i\tanh(G_i)$, and $N^{-1}\sum_i\xi_i^2$ gives a one-step state error at most $C\delta$, except on probability $C\exp(-cN\delta^2)$.

The errors grow only polynomially with $H$. During native transitions the Jacobians are triangular, their diagonal entries have absolute value at most one, and their off-diagonal entries are uniformly bounded. Their products therefore have norm $O(H)$. During the reverse tail let $q_t^*=q_t$ denote the reference variance. The exact scalar states obey $q_\ell\le C_dq_t^*$: the native baseline is $O(1/H)$ and the excess-variance estimate sharpened to the current time gives $\Delta_\ell\le C_ds^2/(t+1)^2$. Also $\alpha_t\le C_dq_t^*$. On enlarged envelopes, the one-step derivative bounds are
\[
 |\partial_q q'|\le1,\quad
 |\partial_r q'|\le C_d(q_t^*)^{3/2},\quad
 |\partial_q r'|\le C_d\sqrt{q_t^*},\quad
 |\partial_r r'|\le1+C_dq_t^*.
\]
The weighted norm $|\Delta q|/\sqrt{q_t^*}+|\Delta r|$ thus grows by at most $1+C_dq_t^*$ per step. The reference variance increases, so changing its denominator at the next step only reduces this norm. Since $\sum_{t=1}^M q_t^*=O_d(\log H)$, the total amplification is polynomial in $H$. The finitely many final steps outside a small-variance interval have bounded derivatives and contribute only a fixed factor. A first-exit argument on enlarged envelopes establishes the same bound for the empirical states.

Take each local accuracy to be a sufficiently small fixed multiple of $\varepsilon H^{-p_0}/\mathcal B$ and also smaller than a fixed multiple of $H^{-1}$, with $p_0$ exceeding the propagation exponent by at least two. A union bound over the $H$ layers, followed by an increase of the fixed polynomial width, makes the good-event readout error at most $\varepsilon/4$ and the bad-event probability at most $c\varepsilon^2/\mathcal B^2$. Give the final readout native covariance and mean $\mathcal B e_H/\sqrt N$. Its mean output is bounded by $\mathcal B$ at every input, and its native conditional variance is at most $1/N$. Consequently a further adjustment of constants gives
\[
 \sup_x\E_\rho|z_w(x)-g(x)|^2\le\varepsilon^2,
\]
using the initial feature approximation and the cubic distortion bound.

For every internal layer the normalized Gaussian relative entropy is exactly $R_\ell^2-1-\log R_\ell^2$. The readout contributes exactly $\mathcal B^2$. The first layer contributes zero. Auxiliary random directions can be included with the same law in the reference probability space; marginalization can only reduce relative entropy. Hence
\[
 \frac2N\KL(\rho\|\mu_{H,N})
 \le\sum_\ell\{R_\ell^2-1-\log R_\ell^2\}+\mathcal B^2
 \le C_d+\mathcal B^2.
\]
This proves Theorem~\ref{thm:continuous-echo}.

\section{A comparison uniform over all smaller variance prefactors}
\label{sec:prefactor-echo}
Let $\nu=Ns\in(0,1]$, so that the original independent Gaussian variances are $\nu/d$, $\nu/N$, and $\nu/N^2$. Put
\[
 \beta=1-\nu,\qquad \kappa(\nu)=1-\nu+\nu\log\nu.
\]
The leading depth cost of the fixed-teacher comparison is $H\kappa(\nu)$, with a bounded remainder at the canonical prefactor.

\begin{theorem}[Uniform comparison through the critical prefactor]
\label{thm:prefactor-echo}
Fix a shallow teacher width, a coefficient bound, an input dimension, and an accuracy $\varepsilon>0$. There are constants $C,A,H_0,K,p$ uniform over that teacher class such that, for every $H\ge H_0$, $N\ge KH^p$, and $0<\nu\le1$, each teacher $g$ admits a comparison law satisfying
\begin{align*}
 2s\KL(\rho\|\mu_{H,N,s})
 &\le (H+d)\kappa(\nu)+C\beta\log(H+1)+A,\\
 \sup_{x\in\mathbb S^{d-1}}\E_\rho|z_w(x)-g(x)|^2&\le\varepsilon^2.
\end{align*}
In particular the normalized cost is at most $C'(1+\beta^2H)$, with $C'$ depending only on the fixed teacher class, dimension, and accuracy. The learner's prior is unchanged. The comparison's first-layer covariance is $I/d$, which agrees with the prior only when $\nu=1$.
\end{theorem}
\begin{proof}
Use the same reverse factors $W_\ell=R_\ell^2$ as in the canonical construction: $W_\ell=1$ during the burn and $W_\ell=1+\alpha_\ell$ during reversal. Inflate the first-layer covariance unconditionally from $\nu I/d$ to $I/d$. Given the current unit direction $e$, use internal row covariance
\[
 \frac1N\{\nu I+(W_\ell-\nu)ee^{\mathsf T}\}.
\]
Its effective variance in direction $e$ is $W_\ell/N$, so the exact scalar recursion is
\[
 q'=F(\nu q+(\beta+\alpha_\ell)r^2),\qquad
 r'=\sqrt{1+\alpha_\ell}\,r(1-q').
\]
The same standardized directional Gaussian defines the next unit direction.

The zero-signal baseline now obeys $b_1=q_0$ and $b_{\ell+1}=F(\nu b_\ell)$. It is bounded above by the canonical native baseline. Define its linear response by
\[
 L_1=1,\qquad
 L_{\ell+1}=\sqrt{1+\alpha_\ell}(1-b_{\ell+1})L_\ell.
\]
Then $c_d\sqrt H\le L_H\le C_dH$, uniformly in $\nu$. The remaining-response bound $L_\ell/L_H\le C_d/(t+1)$ during reversal remains valid: the remaining derivative product is at least its canonical value. During the burn, $L_\ell\le1$.

Start with $r_1=s_0/L_H$, where $s_0$ is the intended final signal. Monotonicity again gives $q\ge b$ and $|r_\ell|\le |s_0|L_\ell/L_H$. The excess variance satisfies the sharper recurrence
\[
 0\le\Delta_{\ell+1}\le\nu\Delta_\ell+
                       (\beta+\alpha_\ell)r_\ell^2.
\]
For $\beta>0$, a disturbance persists for a total weight at most $1/\beta$, and also for at most the remaining number of transitions. The burn contribution to $\sum_\ell\Delta_\ell$ is therefore at most
\[
 \sum_{\ell\text{ in burn}}r_\ell^2
 \le C_ds_0^2 H/L_H^2\le C_ds_0^2.
\]
During reversal, a disturbance with $t$ transitions remaining contributes at most
\[
 C_ds_0^2\min\{t+1,\beta^{-1}\}
 \left(\beta+\frac{C_d}{t+1}\right)\frac1{(t+1)^2}
 \le\frac{C_ds_0^2}{(t+1)^2}.
\]
For $\beta=0$ the canonical proof applies directly. Summing gives $\sum\Delta_\ell\le C_ds_0^2$ uniformly in the prefactor. The correlation-product identity thus gives the same cubic-response bound $|r_H-s_0|\le C_d|s_0|^3$.

Use the same bounded Gaussian-feature encoding of $g_0/(\mathcal B L_H)$ and random unit-direction completion. Because $L_H\ge c_d\sqrt H$, its first-layer concentration bounds and admissibility threshold are uniform in $\nu$. For completeness, finite-width propagation is also uniform. During the burn, the exact states satisfy $q_\ell\le C_d/(\ell+1)$: the excess variance is at most $\max_{j\le\ell}r_j^2\le C_ds_0^2/H$. During reversal, the geometric recurrence gives
\[
 \Delta_\ell\le C_ds_0^2/H+C_ds_0^2/(t+1)^2,
\]
so $q_\ell\le C_dq_t^*$. On either monotone reference-variance segment, write its envelope as $b_*$ and use
\[
 \| (\delta q,\delta r)\|_{b_*,\beta}
 =\frac{K_0\sqrt{b_*}}{\beta+b_*}|\delta q|+|\delta r|.
\]
On small-variance envelopes the derivative matrix is bounded entrywise by
\[
 \begin{pmatrix}
 \nu & C(\beta+b_*)\sqrt{b_*}\\
 C\sqrt{b_*} &1+Cb_*
 \end{pmatrix}.
\]
Choose fixed $K_0$ large enough. The weighted column sums are then at most $1+Cb_*$. Along the burn and reverse segments, $\sum b_*=O_d(\log H)$ and the total variation of the logarithm of the weight is $O_d(\log H)$: its logarithmic derivative with respect to $b_*$ lies between $-1/2$ and $1/2$. The finite number of larger-variance initial or final steps has bounded derivatives uniformly in $\nu$. Hence the total error amplification is polynomial in $H$, uniformly in the prefactor. All actual Gaussian directional variances are bounded by the fixed schedule $W_\ell/N$, so the conditional concentration constants are uniform too. The pointwise mean-square approximation follows exactly as before at $N\ge KH^p$.

The normalized first-layer entropy is $(d+1)\kappa(\nu)$. At an internal transition the exact normalized entropy is
\[
 W_\ell-\nu-\nu\log(W_\ell/\nu)
 =\kappa(\nu)+\alpha_\ell-\nu\log(1+\alpha_\ell)
 \le\kappa(\nu)+\beta\alpha_\ell+\alpha_\ell^2/2.
\]
The readout mean $\mathcal B e_H/\sqrt N$ has normalized cost $\mathcal B^2$, with native covariance $\nu I/N^2$. Since $\sum\alpha_\ell=O_d(\log H)$ and $\sum\alpha_\ell^2=O_d(1)$, summing proves the displayed entropy bound. Finally $\kappa(\nu)\le\beta^2$ follows from its power series, while
$\beta\log(H+1)\le\beta^2H+\log^2(H+1)/(4H)$ and the latter ratio is bounded. This proves the final estimate.
\end{proof}

\section{Logarithmic dependence on the feature-encoding norm}
\label{sec:quantitative-encoding}
The preceding construction requires depth large enough that a small initial signal can be encoded in a unit direction. A short extra amplification stage replaces the quadratic dependence on the encoding norm in that threshold by a logarithmic one.

\begin{theorem}[Quantitative fixed-function comparison]
\label{thm:quantitative-echo}
Fix $d$. Suppose a function on the unit sphere has a bounded Gaussian-feature representation
\[
 g_0(x)=\E[\psi(W)\tanh(W^{\mathsf T}(1,x))],\quad
 W\sim N(0,I/d),\quad \|\psi\|_\infty\le Q,\quad
 \|g_0\|_\infty\le G.
\]
Choose $\mathcal B\ge1+G$ with $C_dG^3/\mathcal B^2\le\varepsilon/4$, and set
\[
 D=\log\{e+C_dQ/\mathcal B\}.
\]
There are constants depending only on $d$ such that, whenever
\[
 H\ge C_dD,\qquad
 N\ge C_d(1+\mathcal B/\varepsilon)^{C_d}H^{C_d}e^{C_dD},
\]
for every $0<\nu\le1$ there is a comparison with pointwise mean-square error at most $\varepsilon^2$ from $g_0$ and
\begin{align*}
 2s\KL(\rho\|\mu_{H,N,s})
 \le {}&(H+d)\kappa(\nu)+\mathcal B^2+C_d\\
 &+C_d\left\{\beta[D+\log(H+1)]+D^2/H\right\},
 \qquad s=\nu/N.
\end{align*}
In particular this cost is at most $C_d[1+\mathcal B^2+\beta^2H+D^2/H]$.
\end{theorem}
\begin{proof}
Put $M=\lfloor(H-1)/3\rfloor$. Use $M$ effective-native burn transitions, then $M$ transitions with effective directional factor
\[
 W_*=e^{2D/M},\qquad \alpha_*=W_*-1,
\]
then the $M$ reverse factors from the preceding construction. Append at most two effective-native transitions. The depth hypothesis makes $\alpha_*\le C_dD/H$ a fixed small constant. The middle block has exact correlation gain $W_*^{M/2}=e^D$ before the tanh derivative factors. The zero-signal baseline still obeys $b'=F(\nu b)$, so
\[
 c_de^D\sqrt H\le L_H\le C_de^DH.
\]
The remaining-response ratio in the reverse block is unchanged. At a state in the middle block with $j$ middle transitions still remaining,
\[
 \frac{L_\ell}{L_H}\le\frac{C_d}{H W_*^{j/2}}.
\]
This follows because the full reverse product is comparable to $H$, and all remaining baseline derivative factors have a product bounded below by a positive constant: the burn has already reduced the baseline variance to $O_d(1/H)$.

The excess-variance argument for the burn and reverse blocks is unchanged. For the middle block, the contribution of its $\beta$ term to the sum of later excess variances is at most $\sum r_\ell^2=O_d(s_0^2/H)$, using the geometric persistence bound when $\beta>0$. Its extra-spike contribution is at most
\[
 C_dH\alpha_*\frac{s_0^2}{H^2}
                 \sum_{j=0}^{M-1}W_*^{-j}
 \le\frac{C_ds_0^2}{H},
\]
because $\alpha_*/(1-W_*^{-1})=W_*$ is bounded. Thus the final distortion remains $|r_H-s_0|\le C_d|s_0|^3$, uniformly in $D,H,\nu$ in the stated range.

The initial direction uses $a_i=\psi(W_i)/(\mathcal B L_H\sqrt N)$. The definition of $D$, after increasing its dimension-dependent coefficient, guarantees $\|a\|\le1/2$. Its feature-average summands are bounded after multiplication by $1/(\mathcal B L_H)$, so their concentration constants can be chosen independent of $Q$. Random orthogonal completion works as before.

The burn and reverse blocks amplify finite-width state errors polynomially in $H$. Throughout the middle block the exact variance is $O_d(1/H)$: its new excess-variance contribution at a given time is at most $C_ds_0^2/H^2$, while the burn contribution is $O_d(s_0^2/H)$. For its derivative estimate one may use the larger envelope
\[
 b_*=(D+1)/H,
\]
so that $\alpha_*\le C_db_*$ and the weighted norm from the prefactor proof gives growth at most $1+C_db_*$ per step. The entire middle block therefore contributes $e^{C_dD}$. Changes of weighted norms at its two boundaries have polynomial cost in $H$; $1\le D\le C_dH$ makes their ratios polynomially bounded. The same first-exit and conditional concentration argument gives the displayed width threshold and pointwise mean-square accuracy.

The middle block adds normalized action at most
\[
 M\{\beta\alpha_*+\alpha_*^2/2\}
 \le C_d(\beta D+D^2/H)
\]
beyond its ordinary depth baseline. First-layer inflation, the other internal layers, and the readout have exactly the costs already calculated. The two appended native transitions add only their baseline. Summing proves the entropy bound. The final simplification follows from $\kappa(\nu)\le\beta^2$, $\beta D\le(\beta^2H+D^2/H)/2$, and the corresponding bound for $\beta\log(H+1)$.
\end{proof}

For a width-$C$ shallow teacher with all coefficients bounded by $B_{\rm teach}$, the Gaussian-feature balls can be chosen so that, at any fixed approximation accuracy,
\[
 \log(e+Q)\le C_d\{1+B_{\rm teach}^2+
                   \log[2+C B_{\rm teach}/\varepsilon]\}.
\]
This follows by using radius proportional to $\varepsilon/(1+C B_{\rm teach})$ and lower-bounding the Gaussian density uniformly on the balls. If the teacher's function is bounded independently of $C$, the readout factor $\mathcal B$ in the comparison can also be chosen independently of $C$. Thus the encoding part of the cost depends only logarithmically on teacher width, and its finite-width requirement is polynomial in that width at fixed dimension, coefficient bound, and accuracy. The theorem still imposes the displayed logarithmic minimum depth; it is not a claim about arbitrary shallow learners at that same sharp cost.

\section{Encoding below the logarithmic depth threshold}
\label{sec:high-gain-encoding}
Large prefix gains also admit a controlled comparison. Their relative entropy is larger, but a total-energy covering argument can use that cost efficiently when depth is at most a constant times the logarithmic encoding norm.

\begin{theorem}[High-gain comparison at short depths]
\label{thm:high-gain-echo}
Use the feature representation, output bound, and readout choice of Theorem~\ref{thm:quantitative-echo}, and put $D=\log(e+C_dQ/\mathcal B)$. For every $H\ge2$ and every prefactor $0<\nu\le1$, a comparison with pointwise mean-square error at most $\varepsilon^2$ exists at every
\[
 N\ge C_d(1+\mathcal B/\varepsilon)^{C_d}H^{C_d}e^{C_dD},
\]
and has normalized entropy
\[
 2s\KL(\rho\|\mu_{H,N,s})
       \le C_d\{1+\mathcal B^2+H\exp(C_dD/H)\}.
\]
The constants are independent of the prefactor. At $H=1$ a direct comparison instead has cost at most $Q^2+(d+1)\kappa(\nu)$ and mean-square error at most $(Q^2+1)/N$ from the represented function.
\end{theorem}
\begin{proof}
First suppose $H\ge4$. Use exactly the three-block schedule from Theorem~\ref{thm:quantitative-echo}, with $M=\lfloor(H-1)/3\rfloor$ and prefix gain $R=e^{D/M}$. We no longer assume $D/H$ is small. The native baseline and its derivative product do not depend on the covariance spikes, so $L_H\ge c_de^D\sqrt H$ still holds. Before a middle-block transition with $j\ge1$ such transitions remaining,
\[
 L_\ell/L_H\le C_d/(H R^j).
\]
The indexing includes the transition about to be performed. Consequently, with $\alpha=R^2-1$,
\[
 \alpha\sum_{j=1}^M R^{-2j}=1-R^{-2M}\le1.
\]
Each middle-block variance disturbance lasts at most $H$ steps, so its extra-spike contribution to the cumulative variance is at most $C_ds_0^2/H$, with no small-gain assumption. The $\beta$ contribution is at most $\sum r_\ell^2\le C_ds_0^2/H$ by geometric persistence. Burn and reverse contributions retain their previous bounds. Thus the same dimension-only cubic response estimate holds, and the same fixed readout factor achieves the requested deterministic approximation.

The entropy contribution of the middle block above its baseline is at most $M\alpha\le H\exp(C_dD/H)$. All other contributions are bounded by $C_d(H+\mathcal B^2)$, which are absorbed by the displayed cost. The native first-feature encoding still fits inside a unit direction because $L_H\ge c_de^D\sqrt H$.

For the finite-width statement, if $H\ge C_dD$ use the polynomial-width proof already given. Otherwise $H\le C_dD$. On the admissible scalar domain $0\le q\le1$, $|r|\le1$, the Jacobian of a middle transition is bounded by $C_dR^3$ in ordinary norm: its entries have magnitudes at most $1$, $2R^2$, $R$, and $3R^3$, respectively. Its $M$-step amplification is at most
\[
 C_d^M R^{3M}\le\exp\{C_d(H+D)\}\le e^{C_dD}.
\]
The burn has polynomial amplification. The reverse block has polynomial amplification on the same envelopes as before; alternatively its global derivative product is bounded by $C_d^M(\prod R_t)^3\le e^{C_dH}H^{C_d}$, which suffices in this short-depth branch. The at most two appended native transitions contribute a bounded factor. Thus all local state errors are amplified by at most $H^{C_d}e^{C_dD}$. Conditional concentration remains uniform even when $R$ is large: hidden squares are bounded, the directional Gaussian $\xi$ is standardized, and $\xi\tanh(G)$ has a universal sub-Gaussian envelope. Enforcing correspondingly small local errors gives the displayed polynomial-in-$e^D$ width. The normalized first-feature summands are bounded because their direction norm is at most one half. Bad-event readout errors are controlled by the fixed factor $\mathcal B$, exactly as before.

For $H=2$ or $3$, use all $H-1$ transitions as a gain block with product $e^D$. Every native derivative product over this fixed number of steps is bounded below by a positive dimension-dependent constant. Hence $L_H\ge c_de^D$, and the remaining-response ratios are bounded by $C_dR^{-j}$. The same geometric identity gives cumulative variance $O_d(s_0^2)$, and therefore cubic distortion. Cost and propagation are bounded by $e^{C_dD}$, which fit the displayed formulas after increasing the dimension-only constant.

Finally let $H=1$. Inflate the first-layer covariance to $I/d$ and conditionally give the readout its native covariance and mean $\psi(W_i)/N$. The mean output is the empirical feature average. Its sampling variance is at most $Q^2/N$, and the native readout adds at most $1/N$. Its normalized readout entropy is $\E N\sum_i[\psi(W_i)/N]^2=\E\psi(W)^2\le Q^2$. The first-layer entropy is $(d+1)\kappa(\nu)$. This proves the last assertion.
\end{proof}

\section{Global state stability and risk comparison without an output cap}
\label{sec:global-response-risk}
A finite-width approximation need not follow a large teacher output outside the label range. The covariance construction gives a same-sign attenuation there. A global stability estimate makes this observation quantitative without any bound on the teacher's output that is independent of its width.

\begin{lemma}[A nonexpansive moment projection]
\label{lem:moment-projection}
Let $\xi,\zeta$ be independent standard Gaussians, let $W\ge\nu>0$, and put
\[
 Y_{r,\sigma}=\tanh(\sqrt W\,r\xi+\sqrt\nu\,\sigma\zeta),
 \qquad
 T_{W,\nu}(r,\sigma)=
 \left(\E\xi Y_{r,\sigma},
   \sqrt{\E Y_{r,\sigma}^2-(\E\xi Y_{r,\sigma})^2}\right).
\]
Then $T_{W,\nu}$ is globally $\sqrt W$-Lipschitz in the Euclidean norm on $(r,\sigma)$, for $\sigma\ge0$.
\end{lemma}
\begin{proof}
Projection in $L^2$ onto the unit vector $\xi$ gives an orthogonal decomposition $Y=r'\xi+Y_\perp$. The map
$Y\mapsto(r',\|Y_\perp\|_2)$ is one-Lipschitz by Pythagoras and the reverse triangle inequality. Also
\[
 \|Y_{r,\sigma}-Y_{\widetilde r,\widetilde\sigma}\|_2^2
 \le W(r-\widetilde r)^2+\nu(\sigma-\widetilde\sigma)^2
 \le W\|(r,\sigma)-(\widetilde r,\widetilde\sigma)\|_2^2.
\]
Compose the two maps.
\end{proof}

\begin{lemma}[Finite-width propagation for arbitrary intended signals]
\label{lem:global-state-propagation}
For each covariance schedule at hidden depth $H\ge2$ used in the preceding comparison constructions, its exact scalar state may be written as $(r,\sigma)$ with $\sigma=\sqrt{q-r^2}$. Suppose the schedule has
$\prod_l\sqrt{W_l}\le C_d(H+1)e^D$, and its initial Gaussian-feature encoding satisfies $\|\psi\|_\infty/(\mathcal B L_H)\le1/2$.
Then, at width
\[
 N\ge C_d(1+\mathcal B/\eta)^{C_d}(H+1)^{C_d}e^{C_dD},
\]
the actual comparison law approximates its deterministic bounded response
$f(x)=\mathcal B r_H(x)$ with
$\sup_x\E_\rho|z_w(x)-f(x)|^2\le\eta^2$.
This holds regardless of the magnitude of the intended signal $g_0(x)/\mathcal B$. The relative-entropy costs of the constructions are unchanged.
\end{lemma}
\begin{proof}
For an empirical transition write $h_i=\tanh(\sqrt W\,r\xi_i+\sqrt\nu\,\sigma\zeta_i)$, where the Gaussian pairs are independent. The next empirical state is
\[
 \widehat r=\frac{N^{-1}\sum_i\xi_i h_i}
                  {\sqrt{N^{-1}\sum_i\xi_i^2}},\qquad
 \widehat q=N^{-1}\sum_i h_i^2,\qquad
 \widehat\sigma=\sqrt{\widehat q-\widehat r^2}.
\]
The radicand is nonnegative by empirical Cauchy--Schwarz. Conditional concentration of the bounded squares, the sub-Gaussian products $\xi_i h_i$, and the chi-square normalization gives errors at most $u^2$ in these three moments except on probability $Ce^{-cNu^4}$. Hence the $(r,\sigma)$ state error is at most $Cu$: use $|\sqrt a-\sqrt b|\le\sqrt{|a-b|}$ only at this single empirical step. The constants are uniform in the state and the factors $W_l$ because $|h_i|\le1$ and $|\xi_i h_i|\le|\xi_i|$.

The initial Gaussian-feature average and random unit-direction completion obey the same state concentration, because their scaled coefficient density is bounded by $1/2$. Apply Lemma~\ref{lem:moment-projection} after each empirical error. Taking
$u=c\eta/[\mathcal B(H+1)^2e^D]$ makes their total propagated readout error at most $\eta/4$. A union bound and the displayed polynomial width make the bad-event probability at most $c\eta^2/\mathcal B^2$. Both the comparison readout mean and the deterministic response have magnitude at most $\mathcal B$ at every input; the native readout noise has conditional variance at most $1/N$. These facts give the pointwise mean-square conclusion.

The next unit direction is still the normalized standardized Gaussian coordinate vector, exactly as in the earlier constructions. Thus this argument changes no comparison law or entropy computation. It replaces only the error-propagation estimate. The fourth power in the local concentration rate changes a polynomial width exponent, rather than iterating a square-root loss through depth.
\end{proof}

At hidden depth $H\ge2$, for every admissible intended signal $s=g_0(x)/\mathcal B$, the scalar comparison response satisfies
\begin{equation}
 0\le\frac{r_H(s)}s\le1\quad(s\ne0),\qquad
 r_H(0)=0,
 \qquad |r_H(s)-s|\le C_d|s|^3.
 \label{eq:response-attenuation}
\end{equation}
The first assertion follows from the exact correlation product: the actual variance is at least its zero-signal baseline, so every ratio of derivative factors lies in $[0,1]$. The cubic estimate follows from the same cumulative excess-variance calculation already given; that calculation bounds the perturbation by $C_ds^2$ and does not require an upper bound on $|s|$. The high-gain schedules have the same estimate because their geometric sum cancels their individual large factors. The encoding guarantees admissibility even for large $g_0$, through its bounded unit-direction norm.

\begin{lemma}[Attenuation is an adequate squared-loss comparison]
\label{lem:attenuation-loss}
Fix $T\ge1$ and $\epsilon>0$. Choose $\mathcal B$, depending only on $d,T,\epsilon$, sufficiently large. For the bounded response $f(x)=\mathcal B r_H(g_0(x)/\mathcal B)$,
\begin{equation}
 (f(x)-y)^2-(g_0(x)-y)^2
 \le\epsilon+y^2\mathbf1_{\{|y|>T\}}
 \label{eq:attenuation-loss}
\end{equation}
for every $x,y$ and every schedule covered by \eqref{eq:response-attenuation}.
\end{lemma}
\begin{proof}
If $|y|\le T$ and $|g_0|\ge2T$, squared loss is no larger at zero than at $g_0$. Convexity therefore makes every same-sign attenuation between those endpoints no worse than $g_0$. If $|g_0|\le2T$, the cubic estimate gives
$|f-g_0|\le C_d(2T)^3/\mathcal B^2$. Choose $\mathcal B$ so large that this changes squared loss by at most $\epsilon$ for $|y|\le T$.
For arbitrary $y$, convexity on the interval from zero to $g_0$ gives
$(f-y)^2-(g_0-y)^2\le\max\{y^2-(g_0-y)^2,0\}\le y^2$.
Combine the cases.
\end{proof}

For bounded labels take $T=1$. For a fixed marginal sub-Gaussian envelope, choose a fixed $T$ so that $\E Y^2\mathbf1_{\{|Y|>T\}}$ is a prescribed small fraction of the accuracy. A single Bernstein event controls the empirical average of this tail variable and the empirical first and second label moments. The loss comparison \eqref{eq:attenuation-loss} is then uniform over all widths and all prefactors, even though the deterministic response itself may vary with the prefactor.

To compare a shallow teacher $g$ with coefficient cap $B_{\rm teach}$ to its Gaussian-feature approximation $g_0$, choose uniform approximation error of order
$\epsilon/(1+CB_{\rm teach}+S)$, where $S$ denotes a fixed label envelope or moment bound. This controls the empirical loss difference on the common label-moment event. The required Gaussian-feature balls have inverse radii polynomial in $C$, so their coefficient density still satisfies
\[
 \log(e+Q)\le C_{d,B_{\rm teach},\epsilon,S}[1+\log(2+C)].
\]
The response readout factor $\mathcal B$ from Lemma~\ref{lem:attenuation-loss} is independent of teacher width. Combining global state propagation with this risk comparison therefore removes the uniform teacher-output-envelope assumption from the polynomial-width oracle.

\section{Uniform shallow-class risk comparisons at every depth}
\label{sec:uniform-risk-comparison}
We now use the bounded response of the covariance construction rather than requiring it to approximate large teacher values. This removes the teacher-output-envelope assumption, while the high-gain schedule removes the logarithmic minimum-depth restriction.

\begin{proposition}[Risk comparison without an output envelope]
\label{prop:uniform-risk-comparison}
Fix $d$, a teacher coefficient cap $B_{\rm teach}\ge1$, and a prescribed accuracy. Put $D_C=1+\log(2+C)$. There are constants $a_0,a_1,a_2$ and a width threshold at most
$(2+C)^{a_0}(H+1)^{a_0}$, uniform over every shallow teacher of width $C$ with that coefficient cap, such that every larger learner width admits a comparison law with the following properties.
For bounded labels, its empirical risk is at most the teacher empirical risk plus the prescribed accuracy, for every deterministic sample. Its normalized relative entropy is at most
\begin{align}
 A_{H,C,\nu}&\le a_1[1+(1-\nu)^2H+D_C^2/H],
     &&H\ge a_0D_C,\label{eq:risk-action-large-depth}\\
 A_{H,C,\nu}&\le a_1H\exp(a_2D_C/H),
     &&1\le H<a_0D_C.\label{eq:risk-action-small-depth}
\end{align}
For a fixed sub-Gaussian label envelope, the same empirical comparison holds on one common label-moment and label-tail event, with constants depending additionally on that envelope. This event is independent of the teacher, learner width, and prefactor.
\end{proposition}
\begin{proof}
We first treat $H\ge2$. Choose a fixed label cutoff $T$: for bounded labels use $T=1$; for a sub-Gaussian envelope choose it so that $\E Y^2\mathbf1_{\{|Y|>T\}}$ is small. Bernstein gives one event on which this empirical tail average is also small and the empirical absolute and squared label moments are bounded by fixed constants. All these choices are independent of teacher width.

Choose the response readout factor $B_0$ using Lemma~\ref{lem:attenuation-loss}; it depends only on dimension, cutoff, and accuracy. Approximate the shallow teacher $g$ by a Gaussian-feature mean $g_0$ with uniform error at most a sufficiently small fixed multiple of
$\epsilon/(1+CB_{\rm teach}+S)$, where $S$ is the fixed label-moment bound. The Gaussian balls may have radius inverse-polynomial in $C$. Hence their coefficient density satisfies
\[
 \log(e+Q)\le c_{d,B_{\rm teach},\epsilon,S}D_C.
\]
The empirical loss difference between $g_0$ and $g$ is small on the same event, since $|g|\le CB_{\rm teach}$ automatically.

Set $D=\log(e+C_dQ/B_0)$. Use the moderate-gain schedule when its depth condition holds, and the high-gain schedule otherwise. The exact response $f=B_0r_H(g_0/B_0)$ is a same-sign attenuation. Lemma~\ref{lem:attenuation-loss} gives its required empirical-risk comparison to $g_0$, uniformly over prefactors. Lemma~\ref{lem:global-state-propagation} approximates this bounded response in pointwise mean square at polynomial width in $H,e^D$. Take its accuracy small enough to control the remaining empirical loss difference using the common label moments. This step does not require $B_0$ to exceed the magnitude of $g_0$: the global state estimate replaces the local stability bounds used in the earlier pointwise-function theorems.

The Gaussian comparison laws and their entropy computations are unchanged. Their normalized actions are those of the moderate-gain and high-gain constructions, with fixed readout factor $B_0$. Since $D\le cD_C$, they give \eqref{eq:risk-action-large-depth} and \eqref{eq:risk-action-small-depth} after increasing constants. Their width thresholds are polynomial in $H,e^D$, hence polynomial in $H,C$.

For $H=1$, draw canonical first-layer Gaussian features and give the readout conditional mean $\psi(W_i)/N$, retaining its native covariance. The normalized entropy is at most $Q^2+(d+1)\kappa(\nu)$, and the pointwise variance around $g_0$ is at most $(Q^2+1)/N$. Choose approximation accuracy inverse-polynomial in $1+CB_{\rm teach}+S$ to transfer squared loss. Both the required width and entropy are at most $\exp(cD_C)$, which has the form claimed at $H=1$. This completes every depth.
\end{proof}

\section{A finite-cover bound for a larger comparison cost}
\label{sec:amgm-large-cost}
For depths below the logarithmic feature-encoding scale, the Sobolev certificate is not the cheapest available estimate. A direct spatial cover retains the arithmetic--geometric mean normalization in the product of layer norms.

\begin{proposition}[AM--GM certificate]
\label{prop:amgm-certificate}
Fix dimension and accuracy, and assume bounded one-Lipschitz hidden activations. For $A\ge1$ choose $\lambda=c_\epsilon/A$. The posterior risk is at most its evidence certificate $F_S(\lambda)$ plus the requested accuracy, from a sample count at most
\begin{equation}
 \exp\left\{C_{d,\epsilon}\left[H+
              (H+1)\log(1+A/H)\right]\right\}
              \operatorname{poly}_{d,\epsilon}(\log(1/\delta)),
 \label{eq:amgm-sample}
\end{equation}
The data event is common to all widths and all variance prefactors if the same bound $A$ is used. In particular, if
$A\le c_1[H\exp(c_2D/H)+B_0^2]$ with fixed $B_0$, the sample count is at most $\exp(C_d(H+D))$ up to fixed-resource factors.
\end{proposition}
\begin{proof}
The total Gaussian operator-energy moment implies
\[
 \Pr_\mu\{Q_{\rm op}>T(H+A+\log(1/\epsilon))\}
 \le\exp[C_0N(H+1)-c_0\omega T(H+A+\log(1/\epsilon))].
\]
For sufficiently large fixed $T$, evidence and the prior risk second moment make the posterior raw-risk contribution of this complement at most a prescribed fraction of $\epsilon$. On the good set put $Q_*=O(H+A)$, suppressing fixed-accuracy constants. Bounded hidden activations give raw amplitude $M\le\sqrt{Q_*}$. The input Lipschitz constant satisfies
\[
 L\le\sqrt{N\|v\|^2}\sqrt{\|W^1\|_F^2/N}
        \prod_{l=2}^H\|W^l\|_\op
 \le\left(\frac{Q_*}{H+1}\right)^{(H+1)/2}.
\]
Use a spatial cover with uniform output accuracy $b=c\epsilon/(M+1)$. Its logarithmic cardinality is at most
$(C_dL/b)^d\log(CM/b)$, with harmless enlargement when $L<1$. Finite-class Hoeffding costs a polynomial in $M,\epsilon^{-1}$ times this log cardinality and a confidence logarithm. Taking logarithms gives a bound of order
\[
 H+(H+1)\log(1+A/H),
\]
which also absorbs $\log(H+A)$ and fixed resource factors. Integrate the good-set deviation, add the posterior tail, and use posterior empirical risk at most $F_S(\lambda)$. A separate empirical-loss comparison can then turn this certificate into an oracle inequality. All dictionaries are independent of width and scale.

If $A\le c_1[H e^{c_2D/H}+B_0^2]$, then
$H\log(1+A/H)\le C(H+D)$, and $H+1\le2H$ for $H\ge1$. This proves the final assertion.
\end{proof}

\section{The smaller-width branch}
\label{sec:small-width-shallow}
The comparison constructions above require a polynomial width threshold. This does not leave smaller realizable widths untreated. A shallow tanh function has a same-width approximate implementation at arbitrary hidden depth, and ordinary finite-parameter covers suffice below the threshold.

\begin{lemma}[A narrow relay implementation]
\label{lem:relay}
Let $g(x)=\sum_{j=1}^C c_j\tanh(a_j^\top\bar x)$ with all coefficients bounded by $B\ge1$, and write $A_0=CB$. For any $H\ge1$ and $0<r<1$, there is a width-$C$, depth-$H$ tanh network $g_H$ with
\[
 \sup_x|g_H(x)-g(x)|\le r.
\]
For $H=1$ take $g_H=g$. For $H\ge2$ it keeps the original first layer, uses a second-layer matrix $\eta I$, uses identity matrices in subsequent hidden layers, and has final readout $c_j/\eta$, where
\[
 \eta=\min\{1,\sqrt{3r/(A_0H)}\}.
\]
In particular its coefficients grow only polynomially in $C,B,H,r^{-1}$.
\end{lemma}
\begin{proof}
The case $H=1$ is immediate. For $H\ge2$, every first-layer feature has magnitude at most one. After multiplying it by $\eta$, every subsequent relay state has magnitude at most $\eta$. The elementary inequality
$|\tanh t-t|\le|t|^3/3$ bounds the accumulated relay error by $(H-1)\eta^3/3$. Dividing by $\eta$ in the readout and summing coefficient magnitudes gives error at most $A_0H\eta^2/3\le r$. If the minimum selects $\eta=1$, its defining inequality gives the same conclusion.
\end{proof}

\begin{lemma}[A comparison at every width]
\label{lem:small-comparison}
Fix input dimension, teacher coefficient cap, and an approximation accuracy. There is a polynomial $P$ in the teacher width and the inverses of the requested accuracies such that, for every $H\ge1$, every $N\ge C$, and every $0<\nu\le1$, a comparison law satisfies
\[
 \int\widehat R_S(z_w)d\rho\le\widehat R_S(g)+\epsilon/32,
 \qquad \frac2\omega\KL(\rho\Vert\mu)\le P(C)H^3.
 \label{eq:small-width-comparison}
\]
The first inequality holds for every deterministic sample on the sphere with labels in $[-1,1]$. Constants in $P$ may depend on the fixed coefficient bound, dimension, and accuracy, but not on depth, learner width, or variance prefactor.
\end{lemma}
\begin{proof}
Use Lemma~\ref{lem:relay} with output accuracy $r$ inverse-polynomial in $1+CB$ and the requested inverse accuracy. At $H=1$, keep the exact teacher and set $\eta=1$ in the following bounds. Clone its $C$ neurons into $N$ nearly equal groups, dividing outgoing coefficients by source-group sizes. The weighted prior energy of this exact implementation is bounded by
\[
 E\le d(d+1)B^2+2(H-1)C+2C\sum_j(c_j/\eta)^2
 \le P_1(C)H.
\]
The second-layer factor $\eta$ only reduces this bound. For $H\ge2$, all internal cloned mean operators have norm at most $\sqrt2$, and the normalized readout norm is at most $\sqrt{2C\sum_j(c_j/\eta)^2}$.

In standard Gaussian coordinates, restrict every first-layer and readout coordinate, and each internal row's $C$ group-constant coordinates, to the interval $[-u,u]$. There are at most $N[d+2+(H-1)C]$ such independent coordinates. The event has probability at least $(3u/8)^{N[d+2+(H-1)C]}$. Conditional on these restrictions, the independent coordinates remain unit-subgaussian; orthogonal unrestricted coordinates remain standard Gaussian. A pair of sphere nets gives conditional probability at least $1/2$ for each internal standardized operator norm to be at most $8\sqrt N$. Independence across layers gives joint conditional probability at least $2^{-(H-1)}$.

The resulting symmetric noise event, translated to the cloned mean, changes each raw output by at most $uP_2(C,H)10^H$, where $P_2$ is polynomial. This follows by the Lipschitz recurrence for normalized hidden errors and the final readout bound. Take $u$ sufficiently small to make this a fixed prescribed output accuracy; then
$\log(1/u)\le O(H)+\log P_3(C,H)$.
Translation of a symmetric Gaussian event costs at most the factor $e^{-\omega E/2}$, by Jensen applied to the remaining Gaussian density ratio. Hence the prior probability of the output neighborhood is at least $\exp[-\omega P(C)H^3/2]$, after enlarging a polynomial $P$ to absorb the logarithms and the fixed dimension.

Condition the prior on this neighborhood. Its normalized entropy is at most $P(C)H^3$. On the sphere, $|g|\le CB$; inverse-polynomial relay and neighborhood accuracies therefore change squared loss by at most a prescribed fraction of $\epsilon$, independently of the sample. This proves \eqref{eq:small-width-comparison}.
\end{proof}

\begin{proposition}[A polynomial-width complement costs polynomially many samples]
\label{prop:small-width-statistics}
Suppose a separate comparison theorem applies above a width threshold $N_0(H,C)$. At every width $C\le N<N_0$, choose
\[
 \lambda_{\rm small}=\frac{c_\epsilon}{P(C)H^3}
\]
with $c_\epsilon$ sufficiently small. There is a common sample event on which these posteriors have risk at most the best shallow comparison risk plus $\epsilon$, using a number of samples polynomial in $H,C,N_0$ and $\log(1/\delta)$. Its constants may depend on the same fixed resources. The event is common to every $0<\nu\le1$.
\end{proposition}
\begin{proof}
The variational identity and Lemma~\ref{lem:small-comparison} bound posterior empirical risk by the comparison empirical risk plus a prescribed fraction of $\epsilon$. A total Gaussian operator-energy cutoff
$Q_{\rm op}\le C[H+\lambda_{\rm small}^{-1}+\log(1/\epsilon)]$
has negligible posterior raw-risk complement: use its exponential moment, the evidence lower bound $e^{-\omega/\lambda}$, and the prior fourth-moment bound on risk. Bounded tanh features also give a raw amplitude bound polynomial in $C,H$ on this good set.

At each $N<N_0$ there are $O(dN+HN^2)$ coordinates with polynomially bounded magnitudes in $N,C,H$. A telescoping estimate bounds the logarithm of a parameter-to-output Lipschitz constant by a polynomial in $H,\log(N+C+H+2)$. Finite parameter grids therefore have log cardinality polynomial in $H,C,N_0$ and the inverse approximation accuracy. Their raw square-loss envelope is polynomial in $C,H$. Finite-class Hoeffding, followed by the uniform output approximation and posterior tail integration, proves the claimed polynomial sample budget. Apply scalar concentration only to one population minimizer in the compact shallow comparison class. All dictionaries are independent of $\nu$.
\end{proof}

\section{Assembly of the all-depth, all-width oracle}
\label{sec:assembly-shallow}
We prove Theorem~\ref{thm:main-shallow}. All constants in this section may depend on the fixed dimension, coefficient cap, accuracy, and confidence. They do not depend on comparison width, learner width, depth, or prefactor.

Proposition~\ref{prop:uniform-risk-comparison} gives a comparison law above a threshold polynomial in $C,H$. Increase the exponent $a$ once so that the theorem's $N_0=(2+C)^a(H+1)^a$ exceeds this threshold and its depth split $H=aD_C$ is valid for the moderate-gain construction. The same choice can dominate the polynomial exponents in the smaller-width comparison below. Choose the three likelihood prefactors sufficiently small to make their variational comparison penalties prescribed fractions of $\epsilon$.

\subsection{Large widths and depths above the encoding scale}
Suppose $N\ge N_0$ and $H\ge aD_C$. Let $j=\lceil\log_2T_{H,C,\nu}\rceil$. Then
$T_{H,C,\nu}\le2^j<2T_{H,C,\nu}$, with a harmless endpoint when $T=1$. The normalized comparison action is at most a constant times $T_{H,C,\nu}$. Thus $\lambda_j=\ell_+2^{-j}$ and the variational identity give, for every sample and every shallow teacher,
\begin{equation}
 F_S(\lambda_j)\le\widehat R_S(g)+\epsilon/16.
 \label{eq:shallow-evidence-comparison}
\end{equation}
The comparison law need not approximate large values of $g$: Proposition~\ref{prop:uniform-risk-comparison} already gives the required squared-loss comparison through attenuation.

Apply the Sobolev certificate with statistical accuracy a fixed small fraction of $\epsilon$. Its energy parameter obeys
$K_j\le c_1 2^j\le c_2T_{H,C,\nu}$, so
\[
 \sqrt{HK_j}\le c_3[\sqrt H+(1-\nu)H+D_C].
\]
Since $H\ge aD_C$, its remaining polynomial sample terms are polynomial in $H$. They are absorbed into an exponential in $\sqrt H$. The factor $e^{cD_C}$ becomes a polynomial in $C$. The resulting budget is therefore at most
\[
 (2+C)^{c_4}\exp\{c_4[\sqrt H+(1-\nu)H]\}.
\]
There are only $O(\log(H+2))$ possible indices $j$, uniformly over prefactors. For a fixed sample count include the indices whose statistical budgets are met and allocate confidence $\delta/[8(J+1)]$ to each. Their union gives one event common to every qualifying prefactor and every learner width. The extra confidence logarithm is absorbed into the same sample bound. On this event posterior risk is at most $F_S(\lambda_j)+\epsilon/8$.

\subsection{Large widths and depths below the encoding scale}
Suppose $N\ge N_0$ and $1\le H<aD_C$. The high-gain comparison, including its direct $H=1$ case, has normalized action at most
$c_5H\exp(c_6D_C/H)$. Increasing $a$ if needed, the theorem's scale
$\lambda=\ell_0[H\exp(aD_C/H)]^{-1}$ makes its variational penalty small. The AM--GM certificate of Proposition~\ref{prop:amgm-certificate} requires at most
$\exp[c_7(H+D_C)]$ samples, after absorbing fixed resources. Because $H<aD_C$, this is a polynomial in $C$. It is dominated by \eqref{eq:main-shallow-budget}. Its fixed dictionary and likelihood scale are common to all prefactors, so no prefactor grid is needed in this branch.

\subsection{All remaining realizable widths}
For $C\le N<N_0$, use Lemma~\ref{lem:small-comparison} and Proposition~\ref{prop:small-width-statistics}. Choose $a$ large enough that
\[
 \ell_-(2+C)^{-a}(H+1)^{-a}P(C)H^3\le\epsilon/64.
\]
At this explicit likelihood scale, the required Gaussian operator cutoff and finite-parameter cover remain polynomial in $C,H,N_0$. Since $N_0$ is polynomial in $C,H$, this sample requirement is polynomial in $C,H$. It is absorbed into \eqref{eq:main-shallow-budget}, and its data event is common to all prefactors. The case $H=1$ uses the exact shallow implementation in the same argument.

The exponent $a$ is chosen first from the comparisons and their thresholds. The sample exponent $c$ is then enlarged to dominate the resulting fixed polynomial powers. This order is noncircular.

\subsection{The leading-one oracle and the sample-count cap}
The bounded coefficient cube is compact, and population risk is continuous. A fixed minimizer $g_*$ in $\mathcal G_{C,B}$ therefore exists. Its automatic output bound is $CB$, so scalar concentration of its squared loss costs only a polynomial in $C$, at fixed coefficient cap and accuracy. Give this event another confidence share. Intersect it with the branch events. The variational comparison, the appropriate statistical certificate, and the smaller-width counterpart yield
$\int R(z)d\Pi\le R(g_*)+\epsilon$, after allocating the error constants. This proves the main theorem simultaneously over its stated widths and prefactors. Jensen gives the ordinary raw posterior-mean bound.

For Corollary~\ref{cor:sample-cap}, the chosen cap $C_n$ is deterministic once sample count, architecture, prefactor, and the fixed theorem resources are specified. Apply the theorem with that cap. No data-dependent model choice or prior mixture is involved.

\section{Sub-Gaussian labels without changing the observation model}
\label{sec:subgaussian-labels}
The bounded-label assumption can be replaced by a fixed sub-Gaussian envelope. This includes a bounded regression function plus independent Gaussian noise. The training labels, likelihood, and prediction rule remain unchanged.

\begin{corollary}[Sub-Gaussian-label oracle]
\label{cor:subgaussian-labels}
Replace $|Y|\le1$ in Theorem~\ref{thm:main-shallow} by
\[
 \E\exp(Y^2/S^2)\le2
\]
for a fixed $S>0$. All its conclusions and sample-complexity exponents remain valid, with constants and likelihood calibration also depending on $S$. No conditional assumption on the label distribution given the input is needed beyond this marginal envelope.
\end{corollary}
\begin{proof}
On one sample event, $n^{-1}\sum_iY_i^2\le cS^2$ and therefore $n^{-1}\sum_i|Y_i|\le cS$, using subexponential Bernstein and a fixed confidence share. Its sample requirement is $O(\log(1/\delta))$ at fixed $S$.
The evidence lower bound becomes
\[
 Z\ge\exp\left[-\frac{\omega}{2\lambda}
                     \left(1+n^{-1}\sum_iY_i^2\right)\right].
\]
Conditioning on hidden features and using the centered Gaussian readout also gives
\[
 \E_\mu R(z)^2\le3+6\E Y^2+\E Y^4\le c(1+S^4).
\]
Thus the same posterior bad-set argument works with
$K=c_S[1+\lambda^{-1}+\log(1/\epsilon)]$ on the common label-moment event.

For a fixed dictionary function with $|f|\le M$, the squared loss $(f(X)-Y)^2$ is subexponential with norm at most $c(M+S)^2$. Finite-class Bernstein therefore replaces finite-class Hoeffding at the same logarithmic-cardinality cost and with the same powers of the output envelope. A uniform output approximation of size $b$ changes empirical or population loss by at most
$2b[M+\E|Y|]+b^2$, with the empirical mean in the empirical version. The label-moment event controls both. Take $b$ of order $\epsilon/(M+S+1)$ in the dictionaries. These changes affect constants and logarithms, not the displayed depth or width exponents.

The comparison laws give pointwise root-mean-square approximation error at most $r$. Their empirical loss exceeds the teacher's by at most
\[
 2r\left(M+n^{-1}\sum_i|Y_i|\right)+r^2,
\]
where $M$ is the envelope of the approximated function. At depths $H\ge2$, the covariance comparison approximates a bounded response whose envelope is fixed independently of $C$. Lemma~\ref{lem:attenuation-loss}, together with one Bernstein event for the fixed label-tail variable $Y^2\mathbf1_{\{|Y|>T\}}$, compares that response to the original teacher. The Gaussian-feature approximation uses inverse-polynomial accuracy in $C$, as in Proposition~\ref{prop:uniform-risk-comparison}. At $H=1$ and in the smaller-width branch, take $M=O(CB)$ and inverse-polynomial $r$. The comparison costs and thresholds remain polynomial in $C$ with the same depth dependence. Finally the fixed teacher's squared loss is subexponential with scale polynomial in $CB+S$, so its scalar concentration event still costs only a polynomial in $C$. Intersect the events and repeat the variational assembly. No clipping of the actual observations or modification of their likelihood has occurred.
\end{proof}

\section{Independent smooth coordinate laws}
\label{sec:nongaussian-law}
Let $P$ have density $p(x)=Z^{-1}e^{-V(x)}$ on $\mathbb R$, mean zero,
and variance one. Assume $V\in C^2$ and
\begin{equation}
 0<m\le V''(x)\le M<\infty\qquad(x\in\mathbb R).
 \label{eq:coordinate-curvature}
\end{equation}
All constants in this section may depend on this one fixed scalar law.
Strong log-concavity implies a sub-Gaussian bound
$\E e^{tX}\le e^{K^2t^2/2}$ for some fixed $K\ge1$.
The assumptions can equivalently be replaced below by this
sub-Gaussian bound, the upper curvature bound, and the stated
centering, variance and integration-by-parts identities.
The law need not be Gaussian or symmetric. For example, a centered,
variance-normalized density proportional to
$\exp\{-x^2/2-a\cos x\}$ with $0<|a|<1/2$ qualifies.

The network prior uses independent copies of $X\sim P$, multiplied by
$\sqrt{\nu/d}$ in the first layer, $\sqrt{\nu/N}$ in internal layers,
and $\sqrt\nu/N$ in the readout. Here $0<\nu\le1$ and
$\omega=N/\nu$. This is one fixed coordinate shape under the same
variance scalings as the Gaussian learner. There is no support
conditioning and no prior over architectures.

\begin{lemma}[Matrix moments with an exact isotropic baseline]
\label{lem:exact-isotropic-moments}
Let $X_1,\ldots,X_N$ be independent vectors with independent
coordinates of law $P$, and let $J$ be a deterministic matrix with
compatible dimensions. Write $\mu=\|J\|_F^2$. There is $c=c_K$ such
that, for every integer $k\ge1$,
\begin{equation}
 \E\left[\frac1N\sum_{i=1}^N\|X_i^\top J\|_2^2\right]^k
 \le \mu^k\prod_{j=0}^{k-1}\left(1+\frac{cj}{N}\right)
 \le\mu^k\left(1+\frac{ck}{N}\right)^k.
 \label{eq:exact-isotropic-moment}
\end{equation}
Consequently, if $W=\sqrt{\nu/N}(X_i^\top)_{i=1}^N$ and $p=2k$,
\begin{equation}
 \E\|WJ\|_F^p
 \le \|J\|_F^p\exp\!\left(\frac{cp^2}{4\omega}\right).
 \label{eq:nongaussian-matrix-moment}
\end{equation}
\end{lemma}
\begin{proof}
The case $\mu=0$ is immediate. Put $Y=\|X_1^\top J\|^2$.
The singular-value decomposition and convexity show that
$\E(Y/\mu)^k$ is at most the largest $2k$-th moment of a unit linear
projection of $X_1$. Such projections are $K$-sub-Gaussian, so
\[
 \E(Y/\mu)^k\le2(2K^2)^k k!\qquad(k\ge1),
 \qquad \E(Y/\mu)=1.
\]
Take $c=16K^4$. For $k\ge2$,
\[
 2(2K^2)^k k!
 \le c^{k-1}(k-1)!
 \le\prod_{j=0}^{k-1}(1+cj).
\]
The first inequality follows from
$32k/8^k\le1$ for $k\ge2$ and $K\ge1$.
Thus every moment of $Y$ is at most the corresponding moment of a
Gamma random variable with shape $1/c$ and scale $c\mu$; the first
moments agree exactly. Expanding the integer moment of a sum and
using independence preserves this termwise domination. The sum of
$N$ independent such Gamma variables has shape $N/c$ and scale
$c\mu$. Its $k$-th moment gives the first inequality in
\eqref{eq:exact-isotropic-moment}; the second is immediate.
Finally,
\[
 \nu^k(1+ck/N)^k
 \le\exp(c\nu k^2/N),
\]
using $t(1+r)\le e^{tr}$ for $0<t\le1$ and $r\ge0$.
This proves \eqref{eq:nongaussian-matrix-moment}.
\end{proof}

\begin{corollary}[The critical Gaussian Sobolev estimate persists]
\label{cor:nongaussian-sobolev}
For a depth-$H$ network with any globally $1$-Lipschitz hidden
activation, its raw readout satisfies, for every even $p\ge2$ and
fixed input $x$,
\begin{equation}
 \E_\pi\|\nabla z_w(x)\|_2^p
 \le\left(\frac{C p}{\omega}\right)^{p/2}
             \exp\!\left(\frac{C H p^2}{\omega}\right).
 \label{eq:nongaussian-gradient}
\end{equation}
For bounded hidden activation, the raw-readout amplitude tail and
prior fourth-risk moment used by the Gaussian evidence certificate
also hold, with constants depending only on $P$ and the label envelope.
Hence the same sample budget
$\exp\{C_d\sqrt{HK}\}$, up to polynomial factors in $H,K$ and
accuracy/confidence, proves the raw posterior evidence certificate
when $K$ is a sufficiently large multiple of $1+\lambda^{-1}$.
\end{corollary}
\begin{proof}
Activation-derivative matrices are contractions and require no
independence. Conditional on earlier layers, apply
Lemma~\ref{lem:exact-isotropic-moments} to each hidden matrix step with the present independent row law.
In the first layer the derivative matrix has $Nd$ independent
coordinates, and the same estimate follows by applying the scalar
sum version of \eqref{eq:exact-isotropic-moment}. Its variance
normalization is $\nu/d$, so the initial Frobenius factor is
$\sqrt N$ and its exponential correction is no larger than the one
in \eqref{eq:nongaussian-matrix-moment}.
Conditional on the last hidden Jacobian, the independent readout is
sub-Gaussian. The same singular-value convexity argument gives its
Hilbert-valued $p$-th moment at most
$(C\sqrt p\sqrt\nu/N)^p$ times the incoming Frobenius norm to power
$p$. Combining the factors proves \eqref{eq:nongaussian-gradient}.

For bounded hidden activation, use
$|z_w(x)|\le\sqrt N\|v\|_2$.
The scalar sub-Gaussian square has a finite exponential moment, so
\[
 \E\exp(c\omega N\|v\|_2^2)\le e^{CN}.
\]
The same conditional readout calculation gives a bounded prior
fourth output moment. The Sobolev proof now uses
$p\asymp\omega\sqrt{K/H}$ exactly as in the Gaussian case; choosing
a sufficiently small proportionality constant absorbs the changed
moment constant. Markov's inequality and Morrey's inequality give
a common Holder-$1/2$ class with Holder constant
$C_d e^{C\sqrt{HK}}$, outside prior mass $e^{-c\omega K}$.
The small-effective-width branch uses ordinary finite-parameter
covers. Evidence lower bounds and Cauchy--Schwarz restore the raw
risk on the exceptional set. Only numerical constants change.
\end{proof}

\begin{lemma}[Rank-one transport and translation entropy]
\label{lem:nongaussian-rank-one-entropy}
Let $X\sim P^{\otimes N}$, let $e$ be a unit vector, and let
$T_a=I+(a-1)ee^\top$ with $a\ge1$. Then
\begin{equation}
 \operatorname{KL}((T_a)_\#P^{\otimes N}\|P^{\otimes N})
 \le a-1-\log a+\tfrac M2(a-1)^2
 \le\tfrac{M+1}2(a-1)^2.
 \label{eq:rank-one-transport-entropy}
\end{equation}
For every deterministic translation $b\in\mathbb R^N$,
\begin{equation}
 \operatorname{KL}(\mathcal L(X+b)\|P^{\otimes N})
 \le\tfrac M2\|b\|_2^2.
 \label{eq:translation-entropy}
\end{equation}
In particular, a physical row with native scale $\sqrt{\nu/N}$
can have its directional variance changed from $\nu/N$ to $W/N$,
for $W\ge1$, at normalized layer cost at most
$(M+1)(\sqrt W-\sqrt\nu)^2$. A readout mean
$\mathcal B e/\sqrt N$ costs at most $M\mathcal B^2$ in normalized
relative entropy $2\operatorname{KL}/\omega$.
\end{lemma}
\begin{proof}
Integration by parts gives
\[
 \E V'(X)=0,\qquad \E[XV'(X)]=1.
\]
The change-of-variables formula for relative entropy gives
\[
 \operatorname{KL}((T_a)_\#P^N\|P^N)
 =\E\sum_i[V((T_aX)_i)-V(X_i)]-\log a.
\]
Taylor's upper bound has first-order expectation $a-1$, because
$\E[(e\cdot X)(e\cdot V'(X))]=1$, and remainder at most
$M(a-1)^2/2$. Since $a-1-\log a\le(a-1)^2/2$ for $a\ge1$, this
proves \eqref{eq:rank-one-transport-entropy}.
Translation has zero first-order expectation and proves
\eqref{eq:translation-entropy}.
For a network layer there are $N$ independent rows and
$a=\sqrt{W/\nu}$. Multiply the row bound by $N$ and normalize by
$2\nu/N$. For the readout, the translation in standardized
coordinates is $\mathcal B\sqrt{N/\nu}\,e$; the same normalization
proves the claim.
\end{proof}

\begin{lemma}[A Gaussian first-layer comparison has bounded action]
\label{lem:nongaussian-first-layer}
Replacing all first-layer coordinates by independent $N(0,1/d)$
coordinates has normalized relative entropy at most a constant
$A_{P,d}$, uniformly in $N$ and $0<\nu\le1$.
\end{lemma}
\begin{proof}
For one coordinate, scaling to the original standardized prior gives
$G/\sqrt\nu$, with $G$ standard Gaussian. Its relative entropy
against $P$ equals
\[
 -\tfrac12\log(2\pi e/\nu)+\log Z+\E V(G/\sqrt\nu).
\]
The upper curvature bound gives
$V(t)\le V(0)+V'(0)t+Mt^2/2$, so this is at most
$A_P+M/(2\nu)+\tfrac12\log\nu$.
There are $N(d+1)$ coordinates; multiply by $2\nu/N$ and use
$\nu\le1$ to obtain a finite $A_{P,d}$.
\end{proof}

\section{Finite-width transfer of the Gaussian covariance construction}
\label{sec:nongaussian-transfer}
Write $\phi=\tanh$. All comparison directions below are determined by
weights and auxiliary randomness, independently of the queried input.
The first layer has the Gaussian comparison law of
Lemma~\ref{lem:nongaussian-first-layer}; only the original prior is
non-Gaussian.

\begin{lemma}[Two-projection Lindeberg estimate]
\label{lem:two-projection-lindeberg}
Let $X=(X_1,\ldots,X_N)$ have independent coordinates of law $P$, let
$e$ be a unit vector, and let $b\in\mathbb R^N$ satisfy
$\|b\|_2\le1$ and $\|b\|_\infty\le N^{-1/2}$.
Put $r=e\cdot b$, $q=\|b\|_2^2$, and, for $W\ge1$ and $0<\nu\le1$,
\[
 a=\sqrt\nu\,b+(\sqrt W-\sqrt\nu)r e,
 \qquad \xi=e\cdot X,\qquad Z=a\cdot X.
\]
Let $(\xi_G,Z_G)$ be the centered Gaussian pair with the same
covariance. If $\theta=\|e\|_\infty$, then
\begin{align}
 |\E\phi(Z)^2-\E\phi(Z_G)^2|
 &\le C_P(1+W)^{3/2}(N^{-1/2}+\theta),\notag\\
 |\E[\xi\phi(Z)]-\E[\xi_G\phi(Z_G)]|
 &\le C_P(1+W)^{3/2}(N^{-1/2}+\theta).
 \label{eq:two-projection-lindeberg}
\end{align}
The covariance identities are exact:
\[
 \E\xi^2=1,\qquad \E Z^2=\nu q+(W-\nu)r^2,\qquad
 \E\xi Z=\sqrt W\,r.
\]
\end{lemma}
\begin{proof}
The covariance formulas use only centering and unit variance.
Moreover $\|a\|_2\le\sqrt W$ and
$\|a\|_\infty\le N^{-1/2}+\sqrt W\,\theta$.
Replace coordinates of $X$ one at a time by standard Gaussians.
The first two Taylor terms cancel because the first two moments
match. For $\phi(a\cdot X)^2$, the third derivative in coordinate
$j$ is bounded by $C|a_j|^3$. The total error is therefore at most
$C_P\sum_j|a_j|^3$.

For $(e\cdot X)\phi(a\cdot X)$, the third derivative is
\[
 (e\cdot X)a_j^3\phi'''(a\cdot X)
       +3e_j a_j^2\phi''(a\cdot X).
\]
During one replacement, separate
$e\cdot X=\xi_{-j}+e_jX_j$. The remaining mixed Gaussian/non-Gaussian
coordinates have covariance identity, so $\E|\xi_{-j}|\le1$.
Their independence from the replaced coordinate, and the fixed third
and fourth moments of the scalar laws, bound the Taylor remainder by
\[
 C_P\bigl(|a_j|^3+|e_j||a_j|^3+|e_j|a_j^2\bigr).
\]
Summing, using $\|a\|_2\le\sqrt W$ and the displayed infinity-norm
bound, proves \eqref{eq:two-projection-lindeberg}. The derivatives of
$t\mapsto\tanh t$ through order three are globally bounded.
\end{proof}

\begin{lemma}[A delocalized finite layer follows the Gaussian state map]
\label{lem:nongaussian-state-layer}
Condition on a previous hidden vector $h(x)$ and a unit direction $e$,
with $|h_i(x)|\le1$. Generate independent standardized rows
$X_i\sim P^N$ and use physical rows
\[
 \sqrt{\nu/N}\,T_{\sqrt{W/\nu}}X_i.
\]
Set $\xi_i=e\cdot X_i$, $e'_i=\xi_i/(\sum_j\xi_j^2)^{1/2}$,
$h'_i(x)=\phi(Z_i)$, and
\[
 q'=N^{-1}\sum_i h'_i(x)^2,\qquad
 r'=N^{-1/2}\sum_i e'_i h'_i(x),\qquad
 \sigma'=\sqrt{q'-(r')^2}.
\]
The zero denominator has probability zero. The associated Gaussian
state map is
\begin{align}
 \mathcal T_W(r,\sigma)
 &=\left(\E[\xi_GY],
      \|Y-\xi_G\E[\xi_GY]\|_{L^2}\right),\notag\\
 Y&=\phi(\sqrt W\,r\xi_G+\sqrt\nu\,\sigma\zeta_G),
 \label{eq:nongaussian-gaussian-state}
\end{align}
where $r=e\cdot h(x)/\sqrt N$, $\sigma^2=N^{-1}\|h(x)\|^2-r^2$,
and $\xi_G,\zeta_G$ are independent standard Gaussians.
For $0<u\le1/4$, except with probability $C e^{-c_PNu^2}$,
\begin{equation}
 \|(r',\sigma')-\mathcal T_W(r,\sigma)\|_2
 \le C_P\left[u+(1+W)^{3/2}
                 (N^{-1/2}+\|e\|_\infty)\right]^{1/2}.
 \label{eq:nongaussian-one-step-state}
\end{equation}
Also, with probability at least $1-\eta$,
\begin{equation}
 \|e'\|_\infty\le C_P\sqrt{\log(2N/\eta)/N},
 \label{eq:nongaussian-delocalization}
\end{equation}
provided $N\ge C_P\log(2/\eta)$.
\end{lemma}
\begin{proof}
Conditionally, the pairs $(\xi_i,Z_i)$ are IID. The variables
$\phi(Z_i)^2$ are bounded, $\xi_i\phi(Z_i)$ have sub-Gaussian tails,
and $\xi_i^2$ are subexponential with mean one. Standard scalar
Chernoff bounds show that their three empirical means differ from
expectation by at most $u$, with failure $C e^{-c_PNu^2}$.
In particular the normalizing denominator squared, divided by $N$,
lies in $[1-u,1+u]$.
Lemma~\ref{lem:two-projection-lindeberg} then controls $q'$ and $r'$
against the Gaussian moments with error
$C_P[u+(1+W)^{3/2}(N^{-1/2}+\|e\|_\infty)]$.
Both moment pairs are positive semidefinite and $|r'|,|r_G|\le1$.
The inequality $|\sqrt a-\sqrt b|\le\sqrt{|a-b|}$ therefore gives
\eqref{eq:nongaussian-one-step-state}.
Finally each $\xi_i$ is $K$-sub-Gaussian. A union bound for their
maximum, together with concentration of $N^{-1}\sum_i\xi_i^2$
around one, proves \eqref{eq:nongaussian-delocalization}.
\end{proof}

\begin{lemma}[Global stability and polynomial finite width]
\label{lem:nongaussian-global-transfer}
The map $\mathcal T_W$ in
\eqref{eq:nongaussian-gaussian-state} is globally $\sqrt W$-Lipschitz.
Suppose $H\ge2$ and a Gaussian covariance construction uses controls
$W_2,\ldots,W_H\ge1$ satisfying
\begin{equation}
 \max_\ell W_\ell\le e^{C_dD},\qquad
 \prod_{\ell=2}^H\sqrt{W_\ell}\le C_d(H+1)e^{C_dD}.
 \label{eq:control-product-bounds}
\end{equation}
Use the same controls with the rank-one transports above, starting
from the Gaussian first-layer comparison and its bounded-density
feature encoder. For every fixed accuracy $\varepsilon>0$ and
readout gain $\mathcal B$, there are constants $A,a$ such that
\begin{equation}
 N\ge A(H+1)^a e^{aD}
 \label{eq:nongaussian-polynomial-width}
\end{equation}
ensures the same bounded deterministic response is approximated in
pointwise mean squared error at most $\varepsilon^2$, uniformly in
the input and $0<\nu\le1$. Constants may depend on $P,d,\varepsilon$
and $\mathcal B$, but not on $H,D,N,\nu$.
\end{lemma}
\begin{proof}
The map from $Y\in L^2$ to
$(\E\xi_GY,\|Y-\xi_G\E\xi_GY\|_2)$ is nonexpansive: orthogonal
projection controls the first coordinate, and the reverse triangle
inequality controls the residual norm. Since $\phi$ is
$1$-Lipschitz and $\nu\le W$, this proves the global Lipschitz claim.

For clarity, the first direction is formed from
$a_i=\psi(W_i^1)/(\mathcal B L_H\sqrt N)$ and an independent uniform
unit vector in $a^\perp$, where the scalar encoder is bounded by
$Q$ and $L_H\ge4Q/\mathcal B$. Then $\|a\|\le1/4$ deterministically;
completing to a unit vector preserves infinity norm
$O(\sqrt{\log(N/\eta)/N})$ with failure at most $\eta$.
The first-layer empirical $q,r$ moments have the same concentration
as in the Gaussian construction. This is a comparison-law operation,
not conditioning the original prior on a support event.

Apply Lemma~\ref{lem:nongaussian-state-layer} at each layer, allocating
failure probability $\eta/H$. Every new direction satisfies the
same infinity-norm bound. Take
$u=C_P\sqrt{\log(2H/\eta)/N}$.
The one-step state error is at most
\[
 C_P(1+\max W_\ell)
       [\log(2NH/\eta)/N]^{1/4}.
\]
Global Lipschitz propagation bounds the total error by this quantity
times $H\prod\sqrt{W_\ell}$, plus the first-layer error with the
same product. The polynomial bounds
\eqref{eq:control-product-bounds} imply that
\eqref{eq:nongaussian-polynomial-width}, with a sufficiently large
fixed exponent, makes the final state error as small as required.
The logarithm of $N$ is absorbed by increasing that exponent.
This analysis is pointwise in the input, with constants uniform over
all inputs; it therefore bounds the supremum of pointwise expected
squared errors without taking an input-space union bound.

The final readout mean is $\mathcal B e_H/\sqrt N$ and its native
centered noise contributes at most $1/N$ to pointwise output MSE.
Both the finite-width mean response and the limiting response have
absolute value at most $\mathcal B$. Choose the total exceptional
probability $\eta$ small relative to $\varepsilon^2/\mathcal B^2$,
and then the on-event state error sufficiently small. This completes
the MSE bound. Auxiliary directions and adaptive transports are
measurable functions of previous weights and independent auxiliary
randomness; entropy is bounded by the conditional chain rule and
convexity under removal of the auxiliaries.
\end{proof}

\section{The all-width oracle for smooth non-Gaussian priors}
\label{sec:nongaussian-oracle}
Let $\mathcal G_{C,B}$ consist of one-hidden-layer tanh networks
\[
 g(x)=\sum_{j=1}^{C'} c_j\tanh(a_j\cdot(1,x)),
 \qquad C'\le C,\qquad
 |c_j|\le B,\qquad\|a_j\|_\infty\le B.
\]
An output bias can be represented by one additional constant hidden
feature; this changes only the width convention by one. No separate
bound on the teacher output is assumed.

\begin{theorem}[A smooth non-Gaussian shallow-comparison oracle]
\label{thm:nongaussian-shallow-oracle}
Fix the coordinate law $P$ satisfying
\eqref{eq:coordinate-curvature}, the input dimension $d$, the teacher
coefficient cap $B\ge1$, and accuracy/confidence
$\varepsilon,\delta\in(0,1)$. Let $X$ have any distribution on the
unit sphere in $\mathbb R^d$. Labels may be bounded by a fixed
constant, or may have a fixed marginal sub-Gaussian envelope.

For every hidden depth $H\ge1$, width $N\ge C$, and $0<\nu\le1$,
use the independent coordinate prior of Section~\ref{sec:nongaussian-law},
tanh hidden activation, a linear raw readout, and the likelihood
\begin{equation}
 d\Pi_{S_n,\lambda}(w)=Z_{S_n,\lambda}^{-1}
 \exp\!\left[-\frac{N}{2\nu\lambda}
                  \widehat R_{S_n}(z_w)\right]d\pi_{N,H,\nu}(w).
 \label{eq:nongaussian-likelihood}
\end{equation}
There are constants $a,c>0$, depending only on the fixed quantities
above, and a deterministic calibration specified below, such that
\begin{equation}
 n\ge(2+C)^c\exp\!\left\{c[\sqrt H+(1-\nu)H]\right\}
 \label{eq:nongaussian-sample-budget}
\end{equation}
implies, with probability at least $1-\delta$,
\begin{equation}
 \int R(z_w)\,d\Pi_{S_n,\lambda}(w)
 \le\inf_{g\in\mathcal G_{C,B}}R(g)+\varepsilon.
 \label{eq:nongaussian-shallow-oracle}
\end{equation}
The risk of the ordinary raw posterior mean is no larger. For fixed
$C,H,n$, one event gives this conclusion for every $N\ge C$ and
every prefactor $\nu$ satisfying \eqref{eq:nongaussian-sample-budget}.

The learner has one prescribed architecture and one independent
coordinate prior. The comparison laws used in the proof do not
modify that prior or introduce a distribution over network sizes.
At $\nu=1$, the depth factor is $\exp(O(\sqrt H))$. The constants and
the degree of the width polynomial may depend on $d,B,\varepsilon$
and the label envelope; the theorem is not a jointly polynomial
bound in these quantities.
\end{theorem}

\paragraph{Calibration.}
Put $D_C=1+\log(2+C)$ and
$N_0=\lceil(2+C)^a(H+1)^a\rceil$.
Choose fixed positive constants $\ell_-,\ell_0,\ell_+$ sufficiently
small. If $N<N_0$, take
$\lambda=\ell_-(2+C)^{-a}(H+1)^{-a}$.
If $N\ge N_0$ and $H<aD_C$, take
$\lambda=\ell_0/[H\exp(aD_C/H)]$.
Otherwise put
\[
 T_{C,H,\nu}=1+(1-\nu)^2H+D_C^2/H,
 \qquad j=\lceil\log_2 T_{C,H,\nu}\rceil,
 \qquad\lambda=\ell_+2^{-j}.
\]
Constants can be enlarged in the order used in the proof. This
calibration depends on the declared comparison class and architecture,
not on unknown teacher coefficients or population risks.

\begin{proof}
We spell out how the Gaussian scalar construction and the preceding
new lemmas yield the claim. This also separates properties of the
comparison law from properties required of the actual prior.

\emph{A bounded comparison response.}
A fixed shallow teacher can be approximated by a bounded Gaussian
feature density:
\[
 g_0(x)=\E[\psi(G)\tanh(G\cdot(1,x))],
 \qquad G\sim N(0,I/d),\qquad \|\psi\|_\infty\le Q.
\]
Place bounded-density Gaussian bumps near its $C$ parameter vectors.
Their radii can be inverse polynomial in $C$ at the required fixed
risk accuracy. Thus $Q\le(2+C)^{a_0}$, uniformly over the class.
In particular $D=\log(e+A_dQ/\mathcal B)=O(D_C)$, where
$\mathcal B$ is a fixed auxiliary readout gain.

The constructions in Theorems~\ref{thm:quantitative-echo} and~\ref{thm:high-gain-echo} use an initial critical
variance-decreasing block, a middle gain block, and a reversed
variance block. Its scalar state is exactly
\eqref{eq:nongaussian-gaussian-state}. It supplies controls satisfying
\eqref{eq:control-product-bounds}, a bounded response
$f_{g,\nu,H}=\mathcal B r_H$, and the action estimates
\begin{align}
 \sum_{\ell=2}^H(\sqrt{W_\ell}-1)^2
 &\le A_d(1+D^2/H)\qquad(H\ge A_dD),\notag\\
 \sum_{\ell=2}^H(\sqrt{W_\ell}-1)^2
 &\le A_d[1+H\exp(A_dD/H)]\qquad(H\ge2).
 \label{eq:nongaussian-schedule-action}
\end{align}
These follow from the critical variance sequence
$q_{j+1}=\E\tanh^2(\sqrt{q_j}G)\asymp(j+1)^{-1}$:
reverse gains telescope, and a gain $e^D$ distributed over a block
of length proportional to $H$ costs $O(D^2/H)$ when $D/H$ is bounded,
or $O(H e^{A_dD/H})$ in general. The schedules for $H=2,3$ use a
fixed number of gain layers and obey the second bound.
The deterministic response has the same sign as $g_0$, has magnitude
at most $|g_0|$, and differs from it by at most
$A_d|g_0|^3/\mathcal B^2$. The scalar construction, including~\eqref{eq:continuous-cubic-response} and~\eqref{eq:response-attenuation}, is unchanged
by the present extension; its ingredients are the Gaussian state map,
the critical variance drift, and the global projection inequality
proved in Lemma~\ref{lem:nongaussian-global-transfer}.

These facts give a risk comparison without a teacher-output cap.
For $|y|\le L$ and $|g_0|\ge2L$, any attenuation of $g_0$ toward
zero cannot increase its squared loss. For $|g_0|\le2L$, choose the
fixed $\mathcal B$ so that the cubic error is sufficiently small.
For all $y$, the excess loss of attenuation is at most $y^2$.
For sub-Gaussian labels, choose fixed $L$ with small
$\E[Y^2\mathbf1_{|Y|>L}]$, and use one common empirical label-tail
and moment event. Taking the Gaussian feature approximation accuracy
inverse polynomial in $C$ controls its loss error as well. Thus the
bounded response has empirical risk at most the original teacher
risk plus the prescribed small error, on a common event.

\emph{A non-Gaussian comparison law.}
Replace only the first-layer comparison distribution by the canonical
Gaussian law. Lemma~\ref{lem:nongaussian-first-layer} charges a fixed
normalized entropy $A_{P,d}$. At subsequent layers use the adaptive
rank-one transports of Lemma~\ref{lem:nongaussian-rank-one-entropy}.
Their normalized action is bounded by
\[
 A_P\sum_\ell(\sqrt{W_\ell}-\sqrt\nu)^2
 \le 2A_P\sum_\ell(\sqrt{W_\ell}-1)^2
                   +2A_PH(1-\sqrt\nu)^2.
\]
Translate the readout mean by $\mathcal B e_H/\sqrt N$ and leave its
centered native noise unchanged. This costs at most $M\mathcal B^2$.
By Lemma~\ref{lem:nongaussian-global-transfer}, width
$N\ge(2+C)^a(H+1)^a$ gives any fixed small pointwise MSE around the
same bounded deterministic response. Consequently the normalized
comparison entropy is at most
\begin{equation}
 A[1+(1-\nu)^2H+D_C^2/H]
 \quad\hbox{if }H\ge aD_C,
 \label{eq:nongaussian-long-action}
\end{equation}
and at most $A[1+H e^{aD_C/H}]$ at smaller $H\ge2$.
The comparison is data-independent. Conditional entropy chain rules
and convexity account for its auxiliary directions.
For $H=1$, use the Gaussian first features and readout mean
$\psi(G_i)/N$. Its normalized entropy and approximation width are
polynomial in $Q$, hence polynomial in $C$; no variance schedule is
needed.

\emph{The large-width statistical bound.}
The variational free-energy inequality bounds the empirical evidence
certificate by the comparison empirical risk plus $\lambda$ times
its normalized entropy. With the stated calibrations this additional
cost is as small as desired. In the long-depth branch,
Corollary~\ref{cor:nongaussian-sobolev} requires sample size, up to
polynomial factors, at most
\[
 \exp\{A_d\sqrt{H[1+(1-\nu)^2H+D_C^2/H]}\}
 \le\exp\{A_d[\sqrt H+(1-\nu)H+D_C]\}.
\]
This has the form \eqref{eq:nongaussian-sample-budget}.
The dyadic entropy levels used in the calibration produce one sample
event for all qualifying $\nu$; their confidence costs are logarithmic.
The function classes in the evidence certificate do not depend on
learner width.

For the short-depth branch use the operator-energy certificate.
Independent sub-Gaussian rows give
\[
 \E\exp(c\omega Q_{\mathrm{op}})\le e^{AN(H+1)},\qquad
 Q_{\mathrm{op}}=N^{-1}\|W^1\|_F^2
     +\sum_{\ell=2}^H\|W^\ell\|_{\mathrm{op}}^2
     +N\|v\|_2^2.
\]
A comparison action $A_*$ therefore permits cutoff
$Q_{\mathrm{op}}=O(H+A_*)$.
The amplitude is $O(\sqrt{H+A_*})$, and AM--GM bounds the input
Lipschitz constant by
$[(H+A_*)/(H+1)]^{(H+1)/2}$.
An input-space cover has logarithmic sample budget
$O_d(H+(H+1)\log(1+A_*/H))$.
For $A_*=O(H e^{aD_C/H})$ and $H<aD_C$ this is $O_d(D_C)$,
so the required sample size is polynomial in $C$. The same reasoning
handles the $H=1$ polynomial action.

\emph{All remaining widths.}
For $C\le N<N_0$, implement the shallow teacher at the same width
and depth by weak identity relays: retain its first hidden features,
use $\eta I$ in the next layer, identity matrices subsequently, and
multiply the readout by $1/\eta$. Taking
$\eta^2=O(\varepsilon/[CB(H+1)])$ makes the accumulated tanh relay
error small. At $H=1$ use $\eta=1$ and the original teacher.
The coefficient bounds and the required parameter-box precision are
polynomial in $C,H,N$; perturbing identity layers by norm
$O(1/H)$ avoids any exponential error amplification.

The curvature upper bound implies $p(t)\ge a_Pe^{-b_Pt^2}$.
For a prior coordinate $sX$ and a target coefficient $\theta$, the
interval $|sX-\theta|\le r$ therefore has negative log probability
at most
\[
 A_P[1+\theta^2/s^2+\log_+(s/r)].
\]
After summing coordinates and multiplying by $2\nu/N$, the box law
has normalized entropy polynomial in $C,H,N$, uniformly in $\nu$;
all coordinate scales satisfy $s\le1$, and $N<N_0$ is itself
polynomial in $C,H$. Choose the small-width likelihood calibration
to pay for this comparison. Ordinary finite-parameter covers have
polynomial sample cost: parameter count is polynomial, and the
logarithm of the parameter-to-function Lipschitz constant is only
$O(H\log(C+H+N))$. Prior operator tails restore raw risk outside the
cover. A countable confidence allocation over widths and norm levels
makes this one event for the entire small-width range.

Finally, concentrate one population-risk minimizing teacher in the
compact coefficient-capped class. Its output bound $CB$ costs only a
polynomial in $C$; for sub-Gaussian labels use the corresponding
subexponential squared-loss concentration. All empirical comparison
steps and the common label-tail event then imply
\eqref{eq:nongaussian-shallow-oracle}. Jensen's inequality gives the
posterior-mean assertion. The constants may be chosen successively:
first comparison accuracies and readout gain, then polynomial width
thresholds, then the small likelihood constants, and finally the
sample exponent $c$.
\end{proof}

\section{Non-log-concave coordinate priors}
\label{sec:gaussian-mixture-priors}
The preceding extension does not fundamentally require log-concavity.
A concrete family with the same guarantee is supplied by finite
Gaussian mixtures with a common positive component variance.

\begin{corollary}[Independent Gaussian-mixture weights]
\label{cor:gaussian-mixture-prior}
Let
\[
 P=\sum_{j=1}^r p_jN(\mu_j,\sigma^2),
 \qquad p_j>0,\quad\sum_jp_j=1,\quad\sigma>0,
\]
with $\sum_jp_j\mu_j=0$ and
$\sigma^2+\sum_jp_j\mu_j^2=1$.
Use independent copies of this one fixed standardized law, scaled by
$\sqrt{\nu/d}$, $\sqrt{\nu/N}$, and $\sqrt\nu/N$ in the three
respective layer types. All conclusions of
Theorem~\ref{thm:nongaussian-shallow-oracle} and its sample-count
calibration hold, with constants allowed to depend on the mixture.
In particular, the sufficient sample count remains
\[
 (2+C)^c\exp\{c[\sqrt H+(1-\nu)H]\}
\]
at every finite $N\ge C$ and every hidden depth $H\ge1$.
The mixture is a scalar coordinate law, not a mixture over network
sizes or architectures. It may be strongly bimodal and need not be
log-concave.
\end{corollary}
\begin{proof}
Write $X=A+\sigma G$, where $A$ takes the values $\mu_j$ with
probabilities $p_j$ and $G$ is an independent standard Gaussian.
Hoeffding's lemma gives a fixed centered sub-Gaussian proxy
\[
 K^2=\sigma^2+(\max_j\mu_j-\min_j\mu_j)^2/4.
\]
Let $p$ be the positive smooth density and put $V=-\log p$.
If $m(x)=\E[A\mid X=x]$, direct differentiation gives
\begin{equation}
 V'(x)=\frac{x-m(x)}{\sigma^2},
 \qquad
 V''(x)=\frac1{\sigma^2}
              -\frac{\operatorname{Var}(A\mid X=x)}{\sigma^4}
 \le\frac1{\sigma^2}.
 \label{eq:mixture-score-curvature}
\end{equation}
The score identities follow either by integration by parts or by
conditioning the latent Gaussian:
\[
 \E V'(X)=\E[G/\sigma]=0,\qquad
 \E[XV'(X)]=\E[(A+\sigma G)G/\sigma]=1.
\]
All required moments are finite. The exact-baseline matrix moments
and the Lindeberg transfer use only isotropy and the fixed
sub-Gaussian proxy. Rank-one transport and translation entropy use
only the upper curvature bound in
\eqref{eq:mixture-score-curvature} and these two score identities;
no positive lower curvature enters those proofs.

For the first-layer Gaussian replacement and the small-width box
comparison, retain any one component to obtain
\[
 p(x)\ge\frac{p_j}{\sqrt{2\pi\sigma^2}}
            \exp[-(x-\mu_j)^2/(2\sigma^2)].
\]
This provides the same quadratic density lower bound and finite
Gaussian-comparison entropy, uniformly after the normalization
$2\nu/N$. Thus every analytic and statistical step of the preceding
proof remains valid, with changed fixed constants. The strictly
positive component variance is fixed; no uniform claim is made as
$\sigma\downarrow0$.
\end{proof}

For example,
$\tfrac12N(\sqrt3/2,1/4)+\tfrac12N(-\sqrt3/2,1/4)$ has mean zero
and variance one, but its negative log density has second derivative
$-8$ at the origin. It is outside the strongly log-concave class and
is covered by Corollary~\ref{cor:gaussian-mixture-prior}.


\section{A polynomial-dimensional Gaussian feature encoder}
\label{sec:dimension-encoder}
The first-layer Gaussian variance decreases with dimension. Approximating a
specified parameter vector by a small Euclidean ball therefore gives a poor
bound on the feature-encoding norm. A finite Hermite expansion avoids that
loss when the comparison rows have a fixed Euclidean norm.

\begin{lemma}[Bounded feature density with polynomial dimension cost]
\label{lem:hermite-encoder}
Fix $C,B\ge1$ and $\eta>0$. There are constants $A,a$, depending only on
these fixed quantities, such that for every $d\ge1$ and every
\[
 g(x)=\sum_{j=1}^C c_j\tanh(a_j^\top(1,x)),\qquad
 |c_j|\le B,\quad\|a_j\|_2\le B,
\]
there is a bounded measurable $\psi$ with
\[
 W\sim N(0,I_{d+1}/d),\qquad
 g_0(x)=\E\psi(W)\tanh(W^\top(1,x)),
\]
\[
 \sup_{x\in\mathbb S^{d-1}}|g_0(x)-g(x)|\le\eta,
 \qquad \|\psi\|_\infty\le A(d+1)^a.
\]
The constants are uniform over the displayed comparison class. The row
norm bound, rather than a coordinatewise bound independent of dimension,
is part of this lemma's hypothesis.
\end{lemma}
\begin{proof}
Choose an odd polynomial
\[P(t)=\sum_{k\le K,\ k\text{ odd}}p_kt^k\]
that approximates $\tanh t$ on $[-\sqrt2B,\sqrt2B]$ within
$\eta/(2CB)$. Such a polynomial follows, for example, by approximating
on this compact interval and taking the odd part. Its degree and
coefficients depend only on the fixed resources.

Let $G\sim N(0,I_{d+1})$, $v=(1,x)/\sqrt2$, and let $H_k$ be the
probabilists' Hermite polynomial. For any unit vector $u$, the Gaussian
regression identity gives
\begin{equation}
 \E H_k(u^\top G)\tanh(\sqrt{2/d}\,v^\top G)
  =b_k(d)(u^\top v)^k,
 \quad b_k(d)=\E H_k(Z)\tanh(\sqrt{2/d}\,Z).
 \label{eq:hermite-ridge-identity}
\end{equation}
This follows by conditioning on $v^\top G$, since
$\E[H_k(u^\top G)\mid v^\top G]=
 (u^\top v)^kH_k(v^\top G)$.
For odd $k$, the coefficient never vanishes and
\begin{equation}
 |b_k(d)|\ge c_kd^{-k/2},\qquad d\ge1,
 \label{eq:hermite-coefficient-lower}
\end{equation}
where $c_k>0$ depends only on $k$. Indeed the Fourier identity
\[
 \tanh t=\int_0^\infty\frac{\sin(ut)}{\sinh(\pi u/2)}\,du
\]
and $\E H_k(Z)e^{itZ}=(it)^ke^{-t^2/2}$ yield, with
$a=\sqrt{2/d}$,
\[
 b_k(d)=(-1)^{(k-1)/2}
       \int_0^\infty\frac{(au)^ke^{-a^2u^2/2}}
                                {\sinh(\pi u/2)}\,du.
\]
Interchanging the integrals is justified by the linear bound for sine
near zero and exponential decay of the denominator at infinity.
Restricting the last positive integral to $1\le u\le2$ proves
\eqref{eq:hermite-coefficient-lower}, since $a\le\sqrt2$.

For $a_j\ne0$ put $u_j=a_j/\|a_j\|_2$. Equations
\eqref{eq:hermite-ridge-identity}--\eqref{eq:hermite-coefficient-lower}
show that the polynomial approximation to $g$ is represented exactly by
\[
 \Psi(G)=\sum_{j,k}c_jp_k
       \frac{(\sqrt2\|a_j\|_2)^k}{b_k(d)}H_k(u_j^\top G).
\]
Terms with $a_j=0$ vanish. Each coefficient in this finite sum has
absolute value at most $A_0(d+1)^{K/2}$.
Replace each $H_k(u_j^\top G)$ by
$H_k(u_j^\top G)\mathbf1_{\{|u_j^\top G|\le R\}}$.
Since $\E H_k(Z)^2=k!$ and
$\Pr(|Z|>R)\le2e^{-R^2/2}$, the uniform error in the represented
function is at most $A_1(d+1)^{K/2}e^{-R^2/4}$.
Take $R^2=A_2+A_3\log(d+1)$ with the fixed constants large enough
that this is at most $\eta/2$. The resulting density is bounded by
\[
 A_4(d+1)^{K/2}(1+R)^K\le A(d+1)^a.
\]
Finally set $\psi(W)$ equal to this truncated density at $G=\sqrt dW$.
Combining the truncation and polynomial approximation errors proves the
claim.
\end{proof}

\section{A dimension-uniform covariance schedule}
\label{sec:dimension-uniform-echo}
The reverse schedule need not be tied to the first-layer variance. Using one absolute reference chain removes dimension dependence from its action and nonlinear response constants.

Let $F(q)=\E\tanh^2(\sqrt q Z)$ and define
\[
 q^{\rm ref}_{-1}=2,\qquad
 q^{\rm ref}_j=F(q^{\rm ref}_{j-1})\quad(j\ge0).
\]
All constants in this section are absolute unless a dependence is displayed. The scalar drift and derivative-product estimates give
\[
 q^{\rm ref}_j\asymp(j+1)^{-1},\quad
 \alpha_t:=\frac{q^{\rm ref}_{t-2}}{q^{\rm ref}_t}-1
       \le\frac C{t+1},\quad
 \sum_{t=1}^M\{\alpha_t-\log(1+\alpha_t)\}\le C,
\]
\[
 \prod_{t=1}^M\sqrt{1+\alpha_t}
 =\left(\frac{q^{\rm ref}_{-1}q^{\rm ref}_0}
              {q^{\rm ref}_{M-1}q^{\rm ref}_M}\right)^{1/2}
 \asymp M+1.
\]
The initialization at two covers every input dimension $d\ge1$.

Fix $0<\nu\le1$, $\beta=1-\nu$, and first-layer comparison covariance $I/d$. Its zero-signal variance sequence is
\[
 b_1=F(2/d),\qquad b_{\ell+1}=F(\nu b_\ell).
\]
By monotonicity, $b_\ell\le q^{\rm ref}_{\ell-1}$ uniformly in $d,\nu$. Therefore
\[
 \prod_{\ell=2}^H(1-b_\ell)\ge cH^{-1/2}.
 \label{eq:dimension-uniform-derivative-product}
\]
After a native burn of at least $H/8$ transitions, the product over any remaining block of at most $H$ such factors is bounded below by an absolute positive constant, because $b_\ell\le C/H$ there.

\begin{lemma}[Uniform response and comparison action]
\label{lem:dimension-uniform-schedule}
Let $D\ge0$ and $H\ge C(D+1)$. Divide the $H-1$ hidden transitions into a native burn of length comparable to $H$, a middle block of length $M\asymp H$, and a reverse block of length $M_r\asymp H$. Use effective directional variances $W=1$ in the burn, $W=e^{2D/M}$ in the middle, and $W_t=1+\alpha_t$ for reverse indices $t=M_r,\ldots,1$.

The linear correlation response $L_H$ obeys
\[
 ce^D\sqrt H\le L_H\le Ce^DH.
\]
If the exact scalar recursion starts at the admissible correlation $r_1=s/L_H$, its final response has the same sign as $s$, satisfies $|r_H|\le|s|$, and obeys
\[
 |r_H-s|\le C|s|^3.
\]
All constants are independent of $d,\nu,H,D$.

Inflating the first-layer covariance from $\nu I/d$ to $I/d$, using these conditional covariance changes, and giving the readout mean $\mathcal B e_H/\sqrt N$ has normalized entropy at most
\[
 C\left[1+\mathcal B^2+\beta^2(H+d)+D^2/H\right].
 \label{eq:dimension-uniform-action}
\]
\end{lemma}
\begin{proof}
The middle directional gain is exactly $e^D$, the reverse gain is comparable to $H$, and the native derivative product has the displayed absolute lower bound. This proves the bounds on $L_H$. If $t$ reverse transitions remain, the remaining directional gain is comparable to $t+1$ and the remaining derivative product is bounded below. Thus
\[
 L_\ell/L_H\le C/(t+1).
\]
If $j$ middle transitions remain, the analogous bound is
$L_\ell/L_H\le C H^{-1}e^{-Dj/M}$. During the burn $L_\ell\le1$.

Write $q_\ell=b_\ell+\Delta_\ell$. The exact recursion gives
\[
 0\le\Delta_{\ell+1}\le\nu\Delta_\ell+(W_\ell-\nu)r_\ell^2,
 \qquad |r_\ell|\le|s|L_\ell/L_H.
\]
A variance disturbance persists for total weight at most the smaller of its remaining number of steps and $1/\beta$, with the latter interpreted as infinity when $\beta=0$. The burn contribution to $\sum\Delta_\ell$ is at most
$Cs^2H/L_H^2\le Cs^2$. For the middle block use
$W-1\le CD/H$ and
\[
 \sum_{j=0}^{M-1}e^{-2Dj/M}\le C M/(1+D).
\]
The $\beta$ source contributes at most $Cs^2/H$, and the $W-1$ source at most $Cs^2/H$. The reverse contribution is bounded by
\[
 Cs^2\sum_{t\ge1}
 \min\{t+1,\beta^{-1}\}
 \left(\beta+\frac C{t+1}\right)\frac1{(t+1)^2}
 \le Cs^2.
\]
Hence $\sum\Delta_\ell\le Cs^2$. The exact product ratio is
\[
 \frac{r_H}{s}=\prod_{\ell=2}^H
      \left(1-\frac{\Delta_\ell}{1-b_\ell}\right)\in[0,1].
\]
Since $1-b_\ell\ge1-F(2)>0$, the elementary product inequality proves the cubic response, with an absolute constant.

The first-layer entropy is $(d+1)\kappa(\nu)$, where
$\kappa(\nu)=1-\nu+\nu\log\nu\le\beta^2$. An internal effective variance $W=1+\alpha$ costs
\[
 W-\nu-\nu\log(W/\nu)
 \le\kappa(\nu)+\beta\alpha+\alpha^2/2.
\]
The middle block contributes at most $C(\beta D+D^2/H)$ beyond its baseline. The reverse block contributes at most $C(1+\beta\log(H+1))$. The readout costs $\mathcal B^2$. Finally
$\beta D\le\beta^2H+D^2/(4H)$ and
$\beta\log(H+1)\le\beta^2H+C$, proving \eqref{eq:dimension-uniform-action}.
\end{proof}

\paragraph{Encoding and finite width.}
Suppose a bounded Gaussian-feature density $\psi$ gives
\[
 g_0(x)=\E[\psi(W)\tanh(W^{\mathsf T}(1,x))],
 \qquad W\sim N(0,I/d),\qquad\|\psi\|_\infty\le Q.
\]
Take $D\ge C+\log(1+Q)$ and $\mathcal B\ge1$. Then
$Q/(\mathcal B L_H)\le1/2$. The usual coefficients
$a_i=\psi(W_i)/(\mathcal B L_H\sqrt N)$ and random orthogonal completion define an exact unit initial direction. Cauchy--Schwarz also gives
$|g_0(x)|/(\mathcal B L_H)\le\sqrt{b_1}/2$, proving initial admissibility even when $b_1\asymp1/d$ is very small. This admissibility does not follow from $L_H\gtrsim\sqrt H$ alone.

All finite-width constants can be dimension-independent. Normalize the state as $(r,\sigma)$ with $\sigma=\sqrt{q-r^2}$. At one transition the random feature is
\[
 Y=\tanh(\sqrt W\,r\xi+\sqrt\nu\,\sigma\zeta),
\]
with independent standard Gaussians. Orthogonal projection in $L^2$ maps $Y$ to
$(\E\xi Y,\|Y-\xi\E\xi Y\|_2)$ nonexpansively. Thus the scalar state update has global Lipschitz constant at most $\sqrt W$. The product of these constants is at most $CHe^D$.

The normalized first-feature summands are bounded independently of $Q,d$, and initial empirical second moments lie in $[0,1]$. Conditional one-step empirical correlations and second moments have absolute Gaussian concentration constants. Passing from second moments to $\sigma$ costs at worst a square root, so an $O(N^{-1/2})$ moment error gives an $O(N^{-1/4})$ state error. The empirical normalization of the directional Gaussian has the same dimension-free control. Union over $H$ transitions and multiply by the global gain bound. For any fixed readout factor $\mathcal B$ and accuracy $\varepsilon$, a sufficient width is
\[
 N\ge C_{\mathcal B,\varepsilon}H^C e^{CD},
\]
with absolute exponents and no hidden dependence on $d$. The readout mean is bounded by $\mathcal B$, and its conditional variance is at most $1/N$, so the bad-event contribution is controlled in mean square as well.

For a bounded target envelope $|g_0|\le U$, one may choose $\mathcal B$ depending only on $U$ and accuracy, by the cubic response estimate. For the bounded-label risk comparison without an output cap on $g_0$, the same-sign attenuation property suffices: if $|Y|\le1$ and $|g_0|\ge2$, every value between zero and $g_0$ has loss at most the loss of $g_0$. On $|g_0|\le2$, a fixed $\mathcal B^2\ge C/\varepsilon$ bounds the cubic distortion by the desired tolerance. These readout choices do not depend on dimension.

\section{Dimension dependence of the gradient certificate}
\label{sec:dimension-sobolev}
A spatial covering argument retains a dimension factor even when the
comparison construction has dimension-independent action. We make that
factor explicit.

\begin{lemma}[A polynomial-dimensional Morrey constant]
\label{lem:dimension-morrey}
For $p\ge2d$, a continuously differentiable function on $[-1,1]^d$
satisfies
\[
 |f(x)-f(y)|\le A d^2\|\nabla f\|_{L^p(\operatorname{Unif}[-1,1]^d)}
                   \|x-y\|_2^{1/2},
\]
where $A$ is absolute. The same assertion holds for its continuous
Sobolev representative.
\end{lemma}
\begin{proof}
Extend $f$ by coordinatewise reflection, with period four. If
$r=\|x-y\|_2\le2\sqrt d$, all balls used below lie in
$[-1-4\sqrt d,1+4\sqrt d]^d$. That box contains at most
$(A_0\sqrt d)^d$ reflected copies of the original cube. Its gradient
$L^p$ norm is therefore at most $A_1d^{1/4}$ times the original
normalized gradient norm, since $p\ge2d$.

Write $v_d$ for the volume of the unit ball. Integrating the fundamental
theorem of calculus along radial segments gives
\[
 |f(x)-\mathop{\rm avg}_{B(x,r)}f|
 \le\frac1{dv_d}\int_{B(x,r)}
              \frac{\|\nabla f(z)\|_2}{\|z-x\|_2^{d-1}}\,dz
 \le2(dv_d)^{-1/p}r^{1-d/p}\|\nabla f\|_{L^p}.
\]
The second inequality follows by polar integration and
$(p-1)/(p-d)\le2$. The averages over $B(x,r)$ and $B(y,r)$ differ
by at most
\[
 r\mathop{\rm avg}_{B(0,r)}\int_0^1
       \|\nabla f(x+z+t(y-x))\|_2\,dt\,dz
 \le v_d^{-1/p}r^{1-d/p}\|\nabla f\|_{L^p}.
\]
Combining these three estimates, using
$v_d^{-1/p}\le A_2d^{1/4}$ for $p\ge2d$, and weakening
$r^{1-d/p}$ to $r^{1/2}$ on $r\le2\sqrt d$ gives the asserted
conservative factor $Ad^2$. Smooth approximation proves the Sobolev
version.
\end{proof}

\begin{proposition}[An explicit-dimensional evidence certificate]
\label{prop:dimension-evidence}
For the Gaussian tanh learner of this paper, fix
$0<\lambda\le1$ and $0<\epsilon,\delta<1$. Put
$K=A(1+\lambda^{-1}+\log(1/\epsilon))$ with a sufficiently large
absolute $A$. One sample event of probability at least $1-\delta$
gives, at every finite width and every variance prefactor,
\[
 \int R(z_w)\,d\Pi\le-\frac{2\lambda}{\omega}\log Z+\epsilon
\]
provided
\begin{equation}
 n\ge\left(\frac{A(d+H+K+1)}{\epsilon\delta}\right)^A
  +\left(\frac{Ad(K+1)}\epsilon\right)^{Ad}
        \exp(A d\sqrt{HK})\log(A/\delta).
 \label{eq:dimension-evidence-budget}
\end{equation}
Increasing the absolute constants permits a fixed geometric grid of
likelihood scales, with the usual summable confidence allocation.
\end{proposition}
\begin{proof}
The Gaussian gradient-moment and readout bounds have absolute constants.
Apply Lemma~\ref{lem:dimension-morrey} in the large-effective-width
case $\omega\sqrt{K/H}\ge d$. Outside prior mass
$2e^{-\omega K}$, the functions have amplitude $O(\sqrt{K+1})$
and half-H\"older constant at most $Ad^2e^{A\sqrt{HK}}$.
An input grid of Euclidean mesh
$[b/(Ad^2e^{A\sqrt{HK}})]^2$ and output quantization at scale $b$
give a fixed uniform cover with
\[
 \log|\mathcal D|\le
       \left(\frac{Ad(K+1)}b\right)^{Ad}
           e^{Ad\sqrt{HK}}\log[A(K+1)/b].
\]
Take $b=c\epsilon/(\sqrt{K+1}+1)$ and apply finite-class Hoeffding
to squared losses. Its polynomial amplitude factor is absorbed into
\eqref{eq:dimension-evidence-budget}.

In the complementary case $N\le\omega<d\sqrt{H/K}$, there are only
polynomially many possible widths. The total operator-energy cutoff
$Q_{\rm op}\le A(H+K)$ has failure mass at most $e^{-\omega K}$,
with absolute constants: the normalized first-layer chi-square
exponent contains $(d+1)/d\le2$. Parameter counts, coordinate caps,
and the logarithm of the parameter-to-output Lipschitz constant are
polynomial in $d,H,K$. Fixed parameter dictionaries therefore require
only the first term of \eqref{eq:dimension-evidence-budget}.
The common raw-risk tail estimate uses $\E_\mu R^2\le24$ and
$Z\ge e^{-\omega/\lambda}$, also dimension-free. It gives the
certificate exactly as in Proposition~\ref{prop:sobolev-oracle}.
\end{proof}

\section{What the upper bound says when dimension grows}
\label{sec:dimension-oracle}
The preceding refinements make a dimension comparison possible without
hiding its entire cost in a constant. The result below fixes comparison
width and bounds the Euclidean norm of each comparison row. It concerns
the Gaussian prior; the earlier non-Gaussian extension was stated with
dimension fixed.

\begin{corollary}[A dimension-explicit shallow oracle]
\label{cor:dimension-shallow-oracle}
Fix comparison width $C$, coefficient bound $B$, accuracy $\epsilon$,
and confidence $\delta$. Bound $|c_j|\le B$ and $\|a_j\|_2\le B$
in the shallow comparison class. There is a constant $A$, depending
only on these fixed resources, such that for every $1\le d\le H$ the
Gaussian learner admits a deterministic likelihood calibration with
\[
 n\ge\exp\{A d[\sqrt H+(1-\nu)H]\}
\]
sufficient for the leading-one comparison inequality
\[
 \int R(z_w)\,d\Pi\le\inf_gR(g)+\epsilon,
 \qquad R(\bar z)\le\inf_gR(g)+\epsilon.
\]
For each $n$, one event of probability at least $1-\delta$ covers all
$N\ge C$ and all qualifying $0<\nu\le1$. Inputs may have any law on
the unit sphere, and labels lie in $[-1,1]$.
\end{corollary}
\begin{proof}
Apply Lemma~\ref{lem:hermite-encoder} at a fixed sufficiently small
accuracy. Its density norm obeys $Q\le A_0(d+1)^{a_0}$ uniformly
over the comparison class. Set $D=A_1\log(d+1)$. The dimension-uniform
schedule gives, for $H\ge A_2\log(d+1)$ and
$N\ge A_3(d+H+1)^{a_3}$, a bounded-response comparison of the desired
risk accuracy and normalized action
\[
 A_4\{1+(1-\nu)^2(H+d)+\log^2(d+1)/H\}.
\]
The original comparison output is automatically bounded by $CB$, so
all response and scalar-concentration constants are fixed here. Put
$T=1+(1-\nu)^2H+\log^2(d+1)/H$ and use
$\lambda=\ell2^{-\lceil\log_2T\rceil}$ for sufficiently small fixed
$\ell>0$. Since $d\le H$, the comparison penalty is a prescribed
small constant. The evidence cutoff satisfies $K\le A_5T$, hence
\[
 \sqrt{HK}\le A_6[\sqrt H+(1-\nu)H+\log(d+1)],
 \qquad K\le A_7(H+1).
\]
Proposition~\ref{prop:dimension-evidence} now gives the asserted
sample bound: the extra $d\log(d+H+1)$ terms are absorbed by
$A d\sqrt H$. A summable confidence allocation over the dyadic
values of $T$ gives one event for all eligible prefactors.

For $C\le N<A_3(d+H+1)^{a_3}$, the same-width weak tanh relay and
Gaussian parameter-box comparison have polynomially bounded
coefficients and inverse accuracy in $d,H$. Their normalized
comparison entropy is polynomial in $d,H$, uniformly in $\nu$.
Use a sufficiently small inverse-polynomial likelihood scale. The
finite-parameter dictionary at the resulting polynomial energy
cutoff has logarithmic cardinality polynomial in $d,H$, with degrees
depending only on the fixed comparison resources. Thus this entire
width branch needs only a polynomial sample count, which is absorbed
by the displayed exponential. The scalar concentration is applied to
one fixed risk minimizer in the compact comparison class, not to a
sample-dependent choice.

Finally, under $d\le H$, the exceptional regime
$H<A_2\log(d+1)$ contains only finitely many integer pairs $(d,H)$,
with bounds depending on $A_2$. The all-depth theorem already proved
handles these pairs at all widths. Increasing $A$ and taking their
finite maximum completes the statement.
\end{proof}

This corollary does not remove the dimension factor from the depth
exponent. For example, at canonical scaling with $d$ of order
$\sqrt H$, it gives $\exp(O(H))$ samples, whereas a lower bound of
$\exp(\Omega(\sqrt H))$ would still leave an exponential gap. More
generally, if $d$ is of order $\sqrt H+(1-\nu)H$, the sufficient
exponent is quadratic in that quantity. Thus a fixed-dimensional
$\exp(O(\sqrt H))$ upper and a varying-dimensional
$\exp(\Omega(\sqrt H))$ lower do not, by themselves, give a matched
depth dependence on continuous domains.


\section{Scope and verification}
The comparison covariance changes only construct a law used in the variational proof. They do not replace or condition the learner's independent prior, whether Gaussian or one of the smooth non-Gaussian laws in the extension. The statistical event is applied to fixed function dictionaries; auxiliary prior-good sets may differ by width and scale. Every learner in the theorem is a finite network.

The proofs have undergone independent AI-agent mathematical review. They have not been formalized in Lean. The accompanying verification record identifies the exact reviewed sources and any qualifications. These results are statistical statements about exact posteriors and do not establish efficient computation.

\end{document}
